\chapter{B 端应用：企业 AI 的千行百业渗透}
\label{ch:enterprise}

%% ============================================================
%% Chapter 10: Enterprise AI — penetrating every industry vertical.
%% From horizontal platforms (Microsoft, Google, Salesforce, ServiceNow)
%% to deep vertical solutions (cybersecurity, supply chain, manufacturing,
%% construction, agriculture, energy, RPA).
%% ============================================================

当消费端 AI 产品在 2023--2025 年间争夺个人用户的注意力时，
另一场规模更加宏大、但远不够``性感''的变革正在全球企业界悄然展开。
从 Fortune 500 的总部大楼到珠三角的工厂车间，
从华尔街的交易大厅到堪萨斯州的农田，
从伦敦的网络安全运营中心到以色列的建筑工地，
AI 正以前所未有的速度渗透进企业运营的每一个毛细血管。

这场渗透并非偶然。当 OpenAI 的 GPT-4 在 2023 年 3 月横空出世时，
全球每一位 CIO 和 CTO 都面临同一个问题：
``这对我们的业务意味着什么？''
答案在过去三年中逐渐清晰——
AI 不是一个可以``观望''的技术趋势，
而是一个正在重新定义每一个行业竞争格局的根本性力量。
率先拥抱 AI 的企业正在建立数据飞轮和运营效率的复利优势，
而观望者则面临被``AI 原住民''竞争对手颠覆的风险。

这场渗透的核心逻辑可以用一个公式来概括：

\begin{keybox}[企业 AI 的 PMF 公式：自动化效率 $\times$ 行业知识深度]
企业 AI 的 Product-Market Fit 取决于两个维度的乘积：
\begin{itemize}[noitemsep]
  \item \textbf{自动化效率}——能否将原本需要人类数小时、数天完成的任务压缩到秒级？
        这是``横向''能力，衡量 AI 对通用工作流的加速程度。
        例如：Microsoft Copilot 将邮件摘要时间从 5 分钟缩短到 5 秒，
        GitHub Copilot 将代码编写速度提升 55\%。
  \item \textbf{行业知识深度}——能否理解特定行业的术语、流程、合规要求和隐性知识？
        这是``纵向''能力，衡量 AI 对垂直领域的真正渗透程度。
        例如：CrowdStrike Charlotte AI 理解 MITRE ATT\&CK 框架的每一个战术和技术；
        Buildots 的 AI 能够区分建筑工地上 200 多种不同的施工状态。
\end{itemize}
纯粹的横向工具（如通用聊天机器人）在企业场景中往往停留在``玩具''阶段——
用户尝鲜后很快回到原有工作习惯；
纯粹的纵向知识（如传统行业软件）则缺乏 AI 带来的效率跃升——
功能完备但使用繁琐，用户培训成本高昂。
只有两者相乘达到阈值，企业客户才会真正买单。

这个公式解释了为什么 Microsoft Copilot 的付费渗透率仅 3.3\%，
而 Glean 却能以 100\% 以上的 NRR 飞速增长——
前者在``自动化效率''维度上做得很好（覆盖面广，功能丰富），
但在``行业知识深度''维度上受限于通用性的天花板；
后者则通过为每个组织构建定制化的知识图谱，
在``行业知识深度''维度上做到了前者难以企及的程度。
\end{keybox}

据多家研究机构估算，2024 年全球企业 AI 市场规模约为 200--240 亿美元。
Verdantix 的研究将 2024 年的市场定价在 130 亿美元
（仅计算企业 AI 平台，不含垂直应用），
并预测以 27.7\% 的 CAGR 增长至 2030 年的 503 亿美元。
Knowledge Sourcing Intelligence 的估算更为激进——
2025 年约 206 亿美元，2030 年将达到 839 亿美元（CAGR 32.4\%）。
Grand View Research 给出了最高的预测：
2030 年市场规模可达 1552 亿美元（CAGR 37.6\%）。
这些数字的巨大差异反映了``企业 AI''定义边界的模糊性——
是仅计算独立的 AI 平台，
还是将嵌入现有企业软件中的 AI 功能也计入？

无论采用哪个数字，结论都是一致的：
企业 AI 是整个 AI 产业中增长最确定、
商业变现最扎实的赛道之一。
与消费 AI 的``免费增长 + 广告变现''模式不同，
企业 AI 从第一天起就有清晰的付费意愿和可量化的 ROI 逻辑。
一个能够将法律合同审查时间从 4 小时缩短到 15 分钟的 AI 工具，
其每用户每月 100 美元的定价几乎不需要论证。

从企业 AI 的发展历史来看，
当前的浪潮实际上是第三波``企业 AI''尝试：

\textbf{第一波（2015--2019 年）：机器学习的企业化}。
以 DataRobot、C3.ai、Palantir 等公司为代表，
核心主张是``将机器学习引入企业决策''。
但这一波的局限在于：
ML 模型需要大量标注数据和专业的数据科学家才能部署，
能力范围局限于预测分析（如客户流失预测、需求预测），
无法处理非结构化数据（文本、图片、音频）。
部署周期长达数月，ROI 难以在短期内量化。

\textbf{第二波（2020--2022 年）：NLP 和对话 AI}。
以 Gong（销售对话分析）、Observe.ai（客服质检）、
Eightfold AI（HR 技能匹配）等公司为代表。
NLP 技术的成熟使 AI 能够理解人类语言，
打开了文本分析、对话分析、情感分析等新的应用场景。
但这些应用仍然是``点状''的——
每个工具只解决一个特定问题，企业需要拼凑数十个工具。

\textbf{第三波（2023 年至今）：生成式 AI 与 Agent}。
GPT-4 和 Claude 等 LLM 的出现彻底改变了企业 AI 的能力边界。
AI 不再局限于``分析''，而是能够``生成''和``行动''——
撰写报告、编写代码、设计演示文稿、处理客户请求、
甚至自主执行复杂的多步骤工作流。
这一波的核心特征是\textbf{通用性}和\textbf{自主性}——
同一个 LLM 可以处理文本、代码、表格、图片等多种模态，
AI Agent 可以在多个系统之间自主导航和执行。

第三波与前两波的根本区别在于：
前两波的企业 AI 需要为每个用例定制一个专用模型，
第三波的企业 AI 可以通过一个通用模型 + 行业知识的组合
覆盖数十甚至数百个用例。
这使得企业 AI 的部署速度从``数月''缩短到``数周''甚至``数天''，
也使得企业 AI 的覆盖范围从``几个核心流程''
扩展到``几乎所有知识工作流程''。

本章将从横向平台到纵向垂直，
系统梳理企业 AI 生态的十个关键战场。

\begin{corebox}[企业 AI 渗透模型：横向平台 vs 纵向垂直]
企业 AI 的竞争格局可以用一个二维矩阵来理解：

\medskip
\noindent
\textbf{横轴}：功能通用性（从``适用于所有行业''到``仅适用于特定行业''）\\
\textbf{纵轴}：AI 深度（从``浅层嵌入——AI 作为现有产品的附加功能''
到``AI 原生架构——产品因 AI 而存在''）

\begin{itemize}[noitemsep]
  \item \textbf{左上象限——横向浅层}：Microsoft Copilot for M365、Google Workspace Gemini。
        覆盖面极广（亿级用户），但 AI 功能以``助手''形态嵌入现有产品，
        渗透深度有限。用户可以完全不使用 AI 功能，产品依然可用。
  \item \textbf{右上象限——横向深层}：Glean、ServiceNow AI、Salesforce Agentforce。
        面向企业通用需求（搜索、IT 服务、CRM），
        但以 AI 为核心重构产品架构，提供 Agent 级别的自主执行能力。
        AI 不是``可选''的附加功能，而是产品体验的核心。
  \item \textbf{左下象限——纵向浅层}：传统行业软件加上 AI 点缀。
        例如在 ERP 系统中添加``智能推荐''功能，
        在 CRM 中添加``AI 生成的邮件建议''。
        本质上并未改变产品逻辑，AI 更多是营销话术而非真实能力。
  \item \textbf{右下象限——纵向深层}：Landing AI（制造质检）、Buildots（建筑监控）、
        Vectra AI（网络安全检测）、Taranis（农业影像分析）。
        从第一天起就为特定行业设计 AI 原生架构，
        拥有该行业独有的数据飞轮和领域知识壁垒。
        没有 AI，这些公司根本不会存在。
\end{itemize}

本章的核心论点是：2024--2026 年企业 AI 的最大趋势，
是从``左上象限''向``右上''和``右下''象限的迁移——
从浅层嵌入走向深层重构，从通用助手走向行业 Agent。
成功的企业 AI 产品要么在``横向深层''做到极致
（如 Glean 的跨平台搜索、Salesforce 的 Agent 化 CRM），
要么在``纵向深层''建立不可替代的行业壁垒
（如 CrowdStrike 的威胁情报、Buildots 的建筑影像数据集）。
停留在``浅层''的产品——无论横向还是纵向——都将面临同质化竞争和价格压力。
\end{corebox}


%% ============================================================
\section{企业 AI 平台}
\label{sec:ent-platforms}
%% ============================================================

企业 AI 平台是这个生态系统中体量最大、影响最广的板块。
四大平台巨头——Microsoft、Google、Salesforce、ServiceNow——
合计服务着全球数十亿企业用户，
它们的 AI 策略将在很大程度上定义
``普通企业员工如何体验 AI''这一基础命题。

\subsection{Microsoft Copilot for M365：巨人的 AI 赌注}

在所有企业 AI 产品中，
Microsoft 365 Copilot 可能是赌注最大、期待最高、
也最具争议的一个。
它是科技史上最大规模的``AI 功能嵌入''尝试——
将 AI 助手同时推送至 4.5 亿付费用户的日常办公工具中。

\subsubsection{产品定位与商业模式}
2023 年 3 月，微软在 GPT-4 发布后仅数周便宣布了 Microsoft 365 Copilot，
将 AI 助手嵌入 Word、Excel、PowerPoint、Outlook、Teams 等核心办公应用。
商业模式直截了当：在现有 Microsoft 365 订阅之上叠加 30 美元/用户/月的 Copilot 附加费。
对于拥有 4.5 亿付费 M365 座席的微软来说，
这意味着一个潜在的 1620 亿美元年化收入增量——
如果每一个现有用户都升级的话。
这个``如果''的分量之重，正是 Copilot 故事的核心张力所在。

Copilot 的技术架构基于三个支柱：
一是 OpenAI 的 GPT-4 系列模型（通过 Azure OpenAI Service 部署）；
二是 Microsoft Graph——微软生态中连接用户数据（邮件、日历、文件、聊天）的 API 层；
三是 Microsoft Search——微软的企业搜索索引。
三者结合，使 Copilot 能够基于用户的个人和组织上下文
生成相关性更高的输出。

\subsubsection{采用现状：1500 万的启示}
然而，截至 2026 年 1 月 28 日（微软 FY2026 Q2 财报），
Copilot 的付费座席数仅为 1500 万，
占 M365 总用户基数的约 3.3\%。
这是微软首次公开披露这一关键指标——
在此前近两年的时间里，微软一直只公布``增长百分比''
而不公布``基数''，这本身就引发了市场的不安。

这一数字首次被公开，在业界引发了复杂的反应。

一方面，1500 万付费用户对应的年化收入约为 54 亿美元
（按 30 美元/月/用户计算），
这在绝对值上已经是一个庞大的数字——
超过了绝大多数独立 AI 公司的全部收入，
也超过了 Slack、Zoom、Notion 等知名 SaaS 公司。
同时，微软透露 Copilot 付费座席同比增长了 160\%，
Copilot 日活用户同比增长了近 3 倍，
显示出强劲的增长动能——
增长率远高于大多数企业软件产品。

另一方面，Forrester 副总裁 J.P.~Gownder 将 3.3\% 的渗透率称为
``令人失望的采用率''（disappointing uptake）。
考虑到微软将整个 M365 的产品策略和市场推广围绕 Copilot 进行了重组，
CEO Satya Nadella 在几乎每一次公开讲话中都将 Copilot 置于核心位置，
这个数字确实低于市场预期。
华尔街分析师此前的预测区间大多在 2000--3000 万付费用户左右。

\subsubsection{渗透率低迷的深层原因}
Copilot 渗透率受限的原因是多维度的，
值得深入分析，因为它反映了企业 AI 采用的普遍性挑战。

\textbf{价格门槛与 ROI 论证困境}。
30 美元/用户/月对大企业来说可能微不足道——
一家 1 万人的企业每月需支出 30 万美元，
年化 360 万美元。对于营收数十亿的大企业，
这只是 IT 预算的零头。
但对中小企业而言，当乘以数百甚至数千个座席时，
就变成了一笔需要严肃论证 ROI 的支出。
而 Copilot 的 ROI 论证恰恰极其困难：
``员工每天节省 15 分钟''如何转化为可量化的财务回报？
是否真的因此减少了加班费？是否因此提高了客户响应速度？
是否因此产生了更高质量的工作成果？
这些问题没有简单的答案，许多 CIO 在试点阶段就卡在了这里。

\textbf{使用场景的模糊性}。
与 GitHub Copilot 在代码补全上的清晰价值不同
（GitHub Copilot 的付费用户超过 190 万，
采纳率在开发者群体中显著更高），
Office Copilot 的使用场景更加分散——
总结会议纪要、起草邮件、制作 PPT、分析 Excel 数据——
每一个场景单独来看都``有点用''，但没有一个是``不可或缺''的。
这与消费 AI 中 ChatGPT 的情况类似：
用户觉得``有用''但很难说清``用了之后我的生产力提高了多少''。
企业用户往往在新鲜感消退后回归原有的工作习惯。

\textbf{数据权限与安全顾虑}。
Copilot 需要访问企业内部的邮件、文档、聊天记录等敏感数据。
许多企业发现，他们现有的数据治理（data governance）并不足以支撑这种
跨应用、跨部门的 AI 数据访问。
Microsoft Graph API 的权限模型是为``搜索''设计的，
不是为``AI 推理''设计的——两者对数据访问粒度的要求截然不同。
一些早期采用者甚至报告了 Copilot ``过度暴露''信息的问题——
员工通过 Copilot 看到了他们本不应该看到的同事邮件或文档内容。
这种``幽灵数据''（shadow data exposure）问题
迫使企业在部署 Copilot 之前先进行大规模的权限清理——
这本身就是一个耗时数月的项目。

\textbf{``浅层嵌入''的局限}。
作为一个覆盖所有行业的通用工具，
Copilot 缺乏对特定行业工作流的深度理解。
一位法律事务所合伙人不需要一个``帮我在 Word 里写更好的段落''的助手，
他需要的是一个理解合同条款、判例法和合规要求的专业 Agent。
一位投行分析师不需要一个``帮我在 Excel 里做更好的图表''的助手，
他需要的是一个能够自主拉取财务数据、构建 DCF 模型、
生成估值报告的金融 Agent。
这些深度行业能力超出了通用 Copilot 的设计范畴，
正是垂直 AI 公司的机会空间。

\textbf{组织变革的惰性}。
最容易被忽略但可能最重要的因素。
引入 AI 助手不仅是一个 IT 决策，更是一个组织变革决策。
它涉及到员工培训、工作流调整、绩效评估标准的修改、
甚至对``什么算是'我的'工作成果''这一基本问题的重新定义。
许多企业的中层管理者对 AI 持抵触态度——
如果 AI 能做我下属 80\% 的工作，
那我还需要管理这么多人吗？我自己的价值又在哪里？
这些``软性阻力''往往比技术障碍更难克服。

\subsubsection{微软的应对策略}
面对渗透率瓶颈，微软在 2025--2026 年采取了多项策略调整：

\begin{itemize}[noitemsep]
  \item \textbf{Copilot Chat（免费层）}：推出免费版 Copilot Chat，
        向所有 M365 订阅用户开放基础 AI 功能，
        通过``先让用户用起来''来培育付费转化的漏斗。
        据 Satya Nadella 透露，Copilot Chat 的用户数是付费 Copilot 的``数倍''。
        这一策略的逻辑是：让用户先体验到 AI 的价值，
        再通过高级功能（如跨应用上下文、Agent 构建）驱动付费转化。
  \item \textbf{Work IQ}：引入基于用户邮件、文件、会议、聊天和组织架构的
        个性化上下文引擎。Work IQ 使 Copilot 的回答更加贴合个人的工作情境——
        不是``通用的 AI 助手''，而是``了解你的日程、你的项目、
        你的同事关系的私人助理''。
        这是微软试图从``横向浅层''向``横向深层''迁移的关键举措。
  \item \textbf{Copilot Studio}：提供低代码/无代码的 Agent 构建平台，
        允许企业定制自己的行业特定 Agent。
        企业可以将内部知识库、行业特定数据和业务规则
        注入 Copilot Studio 构建的 Agent 中，
        从而在``通用平台''的基础上补齐``行业深度''。
        这是微软应对垂直 AI 竞争的核心防线。
  \item \textbf{定价灵活化}：开始针对不同使用场景提供分层定价，
        而非一刀切的 30 美元/月。
        引入基于消费量的定价模式，降低中小企业的采购决策门槛。
  \item \textbf{行业解决方案}：与系统集成商（SI）和独立软件供应商（ISV）
        合作推出行业特定的 Copilot 扩展，
        如金融服务 Copilot、医疗健康 Copilot 等。
\end{itemize}

\subsubsection{战略意义与长期赌注}
无论 Copilot 的短期渗透率如何，
它对微软的战略意义都是不可低估的。
在一个 AI 重塑一切的时代，
微软通过 Copilot 锁定了``企业 AI 的默认入口''这一关键位置。

考虑两个场景：
悲观情景下，Copilot 渗透率在 2028 年达到 10\%
（约 4500 万用户），年化收入约 162 亿美元——
这已经是一个比 Salesforce 全部收入还大的业务。
乐观情景下，渗透率达到 25\%
（约 1.125 亿用户），年化收入约 405 亿美元——
这将使 Copilot 成为微软历史上增长最快的单一产品线。

更重要的是，Copilot 为微软的整个 AI 基础设施
（Azure OpenAI Service、Copilot Studio、Dynamics 365 Copilot、
GitHub Copilot、Security Copilot）
提供了一个统一的企业级分发渠道。
每一个 Copilot 用户都是微软 AI 生态的潜在消费者。

微软的策略本质上是一场``包围战''：
不追求在任何单一垂直领域做到最好，
而是通过无处不在的存在感，
让企业在不知不觉中将 AI 能力嵌入日常工作流。
当 AI 从``选配''变成``标配''时，
微软的 Copilot 将自然而然地成为大多数企业的默认选择——
正如 Windows 成为 PC 的默认操作系统、
Office 成为办公的默认工具一样。
这种``默认效应''（default effect）
是微软在技术变革中反复证明过的核心竞争力。

\subsection{Google Workspace Gemini：后来者的差异化追赶}

Google 在企业 AI 领域的处境颇为微妙。
一方面，Google 拥有全球顶尖的 AI 研究能力（DeepMind、Google Brain），
以及 Gemini 这一强大的多模态模型家族——
在多项基准测试中与 GPT-4 系列不相上下甚至超越。
另一方面，在企业办公软件市场，
Google Workspace 的份额远不及 Microsoft 365。
据估计，Workspace 的企业用户约为 M365 的三分之一至四分之一，
且在大型企业中的渗透率更低——
Fortune 500 公司中使用 M365 的比例超过 90\%。

2026 年 3 月 11 日，Google 宣布了 Gemini AI 在 Workspace 套件中的重大扩展——
这被分析师称为``Google 迄今为止在企业领域最重大的 AI 举措''。
核心升级包括：

\textbf{跨应用上下文感知}（Contextual Awareness）。
Gemini 现在能够跨越 Drive、Gmail、Docs、Sheets 等应用，
从用户的整个数字历史中拉取相关数据。
例如，在 Docs 中起草项目提案时，
用户可以要求 Gemini 基于上个月的特定电子表格制定预算，
并整合最近邮件中的反馈意见。
这种跨应用的信息整合能力是 Google 的差异化优势之一——
因为 Google 本身就是一家``搜索''公司，
跨数据源的检索和关联是其核心基因。
微软虽然有 Microsoft Graph，但 Google 在信息检索领域
拥有二十多年的技术积累。

\textbf{``Help me create''功能}。
在 Docs 中，Gemini 能够根据用户指令创建完整的文档草稿，
自动从其他 Workspace 文件中拉取上下文信息。
在 Sheets 中，Gemini 可以自动填充电子表格、生成公式和图表。
在 Slides 中，AI 可以从文档自动生成演示文稿。
这些功能的设计理念是``消除空白页问题''——
创作的最大障碍往往不是``写''，而是``开始写''。

\textbf{写作风格匹配}（Match Doc）。
``Match doc''功能可以让 Gemini 在添加新内容时
自动匹配现有文档的写作风格和格式。
这在多人协作的场景中尤为实用——
确保来自不同作者的内容风格一致，
减少编辑和格式化的时间。

Info-Tech 研究集团首席研究总监 Julie Geller 评价道：
``这些更多是渐进式改进而非革命性功能，
但它们解决了真实的工作流痛点。
真正的价值在于，Google 正在将 AI 辅助直接嵌入
人们每天使用的工具中——如 Docs、Sheets 和 Slides。''

\subsubsection{Google vs Microsoft：企业 AI 的双寡头对抗}
在企业 AI 平台层面，Google 和 Microsoft 的竞争呈现出一种``不对称对抗''的格局：

\begin{itemize}[noitemsep]
  \item Microsoft 的优势在于\textbf{企业关系深度}——
        它与全球大多数 Fortune 500 公司有数十年的销售关系和技术锁定，
        Active Directory/Entra ID 是企业身份认证的事实标准，
        Azure 是许多企业的首选云平台，
        SharePoint/Teams 是企业协作的核心基础设施。
        这些``锚点''使得 M365 Copilot 的分发成本极低。
  \item Google 的优势在于\textbf{AI 模型能力}和\textbf{搜索基因}——
        Gemini 模型家族在多模态能力上有独特优势
        （长上下文处理、视频理解、跨模态推理），
        Google 在信息检索和知识图谱方面的积累是 Microsoft 所不具备的。
        此外，Google 的 Android 和 Chrome 生态为其提供了
        微软所没有的移动端和浏览器端触点。
  \item 两者都面临同样的挑战：
        如何将``通用 AI 助手''转化为企业用户愿意持续付费的``不可或缺的工具''。
        目前两者都尚未完全破解这一难题。
\end{itemize}

\subsection{Salesforce Agentforce：从 Einstein 到 Agent 的华丽跃迁}

Salesforce 是企业 AI 赛道中最引人注目的``逆袭''故事之一。
其发展轨迹几乎完美地展示了
``从 Copilot 到 Agent''的范式迁移如何重塑一家企业软件巨头的命运。

\subsubsection{Einstein GPT 阶段（2023--2024）：防御性的 AI 包装}
2023 年 3 月，Salesforce 推出 Einstein GPT，
将生成式 AI 嵌入其 CRM 平台。
功能包括：AI 生成的销售邮件、客户洞察摘要、
自动化的客户服务回复建议等。
这在当时更多被视为一种防御性举措——
``竞争对手都有 AI 了，我们也得有。''

市场反应平平，甚至出现了恐慌情绪。
华尔街开始讨论``SaaSpocalypse''——
如果 AI 能够替代人类的销售、客服、营销工作，
那么 Salesforce 基于``按人头收费''的商业模式就面临根本性威胁。
CRM 系统管理的是``人''的活动——
如果活动的执行者从人变成了 AI Agent，
``每个人 25 美元/月''的定价逻辑还成立吗？

\subsubsection{Agentforce 阶段（2024 年底至今）：攻击性的 Agent 重构}
2024 年底，Salesforce 推出了 Agentforce，
这标志着从``Copilot 范式''（AI 作为人类助手）
到``Agent 范式''（AI 作为自主执行者）的根本转变。
Agentforce 允许企业部署自主 AI Agent，
这些 Agent 能够独立处理客户服务请求、执行销售跟进、
管理营销活动、生成分析报告等任务，而无需人类的逐步指导。

关键的战略决策在于\textbf{定价模式}。
Agentforce 不按``用户数''收费，而是按``对话数''收费——
每次 AI Agent 与客户的交互计费。
这一定价创新一举化解了``SaaSpocalypse''叙事：
AI Agent 不是在蚕食 Salesforce 的收入，
而是在创造一个全新的收入流。
一个企业可能只有 100 名人类销售，
但它的 AI Agent 可以同时处理 10000 个客户对话——
按对话计费的收入远超按人头计费。

\subsubsection{114\% 的增长奇迹}
结果令人震惊。截至 2026 年 3 月，
Salesforce 报告其 AI 平台的年化经常性收入（ARR）
实现了 114\% 的同比增长。
这一增速远超市场预期，
也让``SaaSpocalypse''的叙事彻底翻转为``Agentic Renaissance''。
Salesforce 的股价在财报发布后显著上涨，
市场终于相信：对于 B2B SaaS 巨头来说，
AI 不是威胁，而是增长引擎——
\textbf{前提是你能从``嵌入 AI 功能''升级到``以 AI Agent 为核心重构产品''。}

Einstein GPT 阶段的 Salesforce 是在原有产品上``贴''了一层 AI——
CRM 的核心逻辑不变，AI 只是让某些操作更方便。
Agentforce 阶段的 Salesforce 是让 AI 成为了产品体验的中心——
CRM 不再只是``记录客户数据的工具''，
而是``自主管理客户关系的智能平台''。
这两种策略的商业效果天差地别。

Salesforce 创始人兼 CEO Marc Benioff 将 Agentforce 称为
``Salesforce 历史上最重要的产品发布''，
并预言企业软件的未来将从``Software as a Service''
演变为``Agents as a Service''。

\subsection{ServiceNow AI：IT 服务管理的 Agent 化}

ServiceNow 是另一个成功拥抱 AI 的企业软件巨头，
且其策略比 Salesforce 更加激进——
它不仅在自己的产品中嵌入 AI，还通过收购来加速 AI 能力建设。

\subsubsection{Now Assist 的崛起}
ServiceNow 的核心业务是 IT 服务管理（ITSM）和企业工作流自动化。
2023 年，公司推出了 Now Assist——一个基于生成式 AI 的助手，
嵌入到 IT 服务、客户服务、HR 服务等模块中。
Now Assist 的核心功能包括：
自动生成工单摘要、智能分类和路由、
AI 驱动的知识库搜索、以及自然语言驱动的工作流创建。

到 2025 年底，Now Assist 的表现超出预期。
CFO Gina Mastantuono 在 2025 年 5 月的年度大会（Knowledge 2025）上宣布，
Now Assist 的年化合同价值（ACV）已超过 2.5 亿美元，
并预计到 2026 年将达到 10 亿美元。
一些客户在升级 AI 功能后，对 ServiceNow 产品的支出增加了 60\%。
这个数据揭示了一个重要的商业逻辑：
\textbf{AI 功能不是在稀释现有 SaaS 的 ARPU，
而是在大幅提升 ARPU}——
客户愿意为 AI 功能支付显著的溢价，
因为它切实地减少了 IT 运维人力成本。

\subsubsection{28.5 亿美元收购 Moveworks}
2025 年 3 月，ServiceNow 以 28.5 亿美元收购了 Moveworks——
这是 ServiceNow 历史上最大的一笔收购，
也是 2025 年企业 AI 赛道最重大的并购事件之一。

Moveworks 是一家专注于企业 AI 助手的公司，
成立于 2016 年，由 Bhavin Shah、Vaibhav Nivargi、
Varun Singh 和 Jiang Chen 联合创立。
公司此前融资超过 3 亿美元（包括 Tiger Global、
Alkeon Capital、Lightspeed Venture Partners、
Kleiner Perkins 等明星投资方），
2021 年估值达到 21 亿美元。

Moveworks 的核心能力包括两个方面：

\textbf{前端 AI 助手}——Moveworks 的对话式 AI 界面
能够理解和解决企业 IT 问题。
员工可以用自然语言描述问题
（如``我的 VPN 无法连接''或``我需要申请一台新笔记本电脑''），
Moveworks 自动理解意图、查询知识库、
执行操作或创建工单。

\textbf{企业搜索技术}——跨多个企业应用的统一搜索，
能够在 Slack、Confluence、SharePoint、ServiceNow 等
数十个数据源中检索和整合信息。

收购 Moveworks 使 ServiceNow 获得了强大的前端能力
（Moveworks 的对话界面 + 企业搜索），
加上其自身的后端优势（工作流自动化 + ITSM 平台），
形成了``AI 前端 + 自动化后端''的完整闭环。
此外，Moveworks 带来了 350 多家大型企业客户，
包括 Broadcom、Palo Alto Networks、Pinterest 等，
以及 500 多万平台用户。

ServiceNow CFO 确认这次收购``不是对有机增长的替代''，
仅 Moveworks 一项就预计将提升 100 个基点的收入增长。

\subsubsection{ServiceNow 2025 全年财务表现}
ServiceNow 的 2025 全年财报同样亮眼：
Q4 订阅收入 34.66 亿美元，同比增长 21\%（恒定汇率下 19.5\%）；
当前剩余履约义务（cRPO）128.5 亿美元，同比增长 25\%；
剩余履约义务总额 282 亿美元，同比增长 26.5\%；
调整后运营利润率超过 31\%，自由现金流增长 34\%，
自由现金流利润率达到 34.5\%。

CEO Bill McDermott 将 ServiceNow 定位为
``企业业务重塑的 AI 控制塔''（AI Control Tower for Business Reinvention）。
他在财报电话会上表示：``ServiceNow 显著超出了 Q4 预期，
加速了新业务增长，并发布了 2026 年的出色指引。
凭借我们一贯的 Rule of 55+ 表现，
没有任何一家企业 AI 公司比 ServiceNow 更有能力
实现可持续的盈利性收入增长。''

2026 年的指引同样令人印象深刻：
订阅收入预计在 155.3--155.7 亿美元之间
（意味着约 19--20\% 的同比增长）。
Now Assist 的净新增 ACV 在 Q4 2025 同比翻了一倍以上。
ServiceNow 还宣布了 50 亿美元的股票回购计划，
并计划立即执行 20 亿美元的加速回购。

\subsubsection{横向平台的共同困境与机遇}

综合来看，四大企业 AI 平台呈现出几个共同特征：

\textbf{从 Copilot 到 Agent 的范式迁移}。
2024 年的关键词是``Copilot''（副驾驶），
2025--2026 年的关键词是``Agent''（自主代理）。
Salesforce Agentforce、ServiceNow Agent、
Microsoft Copilot Studio、Google Gemini Agents——
所有平台都在从``辅助人类做决策''转向``替代人类做执行''。
这一转变的驱动力来自两方面：
一是 LLM 推理能力的提升使 Agent 的可靠性达到了商业可用的阈值；
二是企业对``自动化''的期望从``提效''升级到``替代''——
不是让员工更快地做同样的事，而是让 AI 直接做员工的事。

\textbf{数据壁垒即护城河}。
每个平台的竞争力最终取决于它能访问多少企业内部数据。
Microsoft 拥有 Outlook/Teams/SharePoint 的数据图谱，
Salesforce 拥有 CRM 客户互动数据，
ServiceNow 拥有 IT 工单和工作流数据，
Google 拥有 Gmail/Drive/Calendar 的数据。
这些数据资产是 AI 功能差异化的根基——
AI 模型的基础能力可以通过调用 OpenAI/Anthropic 的 API 获得，
但\textbf{对特定组织数据的深度理解}是无法外包的。
这也解释了为什么 ServiceNow 愿意花 28.5 亿美元收购 Moveworks——
它买的不仅是技术，更是 Moveworks 在 350 多家企业中
积累的 AI 助手交互数据和企业知识图谱。

\textbf{定价模式的重构}。
按人头（per-seat）收费的传统 SaaS 模式正在被
按使用量（consumption-based）或按结果（outcome-based）收费所补充。
Salesforce Agentforce 按``对话数''而非``用户数''定价；
Microsoft 为 Copilot Studio 引入了消费量定价；
ServiceNow 的 AI 定价与平台使用量挂钩。
这一趋势的本质是：
当 AI Agent 能够``做''人类的工作时，
``人''不再是定价的锚点，``工作量''或``价值''才是。


%% ============================================================
\section{企业搜索与知识管理}
\label{sec:ent-search}
%% ============================================================

如果说办公套件（M365、Workspace）是企业 AI 的``通用入口''，
那么企业搜索则是``垂直深度''的最佳体现。
一个企业平均使用 200 到 400 个 SaaS 应用
（Slack、Confluence、JIRA、Salesforce、Google Drive、
SharePoint、ServiceNow、GitHub、Notion、Asana、HubSpot\ldots），
知识散落在这些孤岛中。
麦肯锡的研究显示，知识工作者每天花费约 1.8 小时
（占工作时间的 20\% 以上）在``寻找和整理信息''上。
AI 驱动的企业搜索正在从根本上解决这一痛点。

\subsection{Glean：企业搜索赛道的绝对王者}

在企业 AI 搜索领域，Glean 是当之无愧的标杆企业。
这家由前 Google 搜索架构师 Arvind Jain 在 2019 年创立的公司，
在短短几年内实现了惊人的增长，
成为 2024--2026 年间最受关注的企业 AI 创业公司之一——
没有之一。

\subsubsection{创始人与产品基因}
Arvind Jain 在 Google 工作了近 20 年，
是 Google 搜索基础设施的核心架构师之一。
他参与构建了 Google 内部的知识检索系统——
一个为超过 18 万 Google 员工提供跨系统信息检索的平台。
这段经历让他深刻理解了两件事：
一是大型组织中``信息碎片化''问题的严重程度；
二是用 AI 解决这个问题的巨大商业价值。

2019 年，Jain 离开 Google，创立了 Glean（最初名为 Scio Technologies）。
创始团队几乎清一色来自 Google，
将 Google 级别的搜索技术与企业知识管理需求结合，
创造了 Glean 的核心产品。

\subsubsection{技术架构三层模型}
Glean 的技术架构基于三个核心层次：

\textbf{连接层}（Connectors）。
Glean 能够与 100 多个企业应用进行深度集成，
包括 Slack、Confluence、JIRA、Google Workspace、Microsoft 365、
Salesforce、ServiceNow、GitHub、Notion、Zendesk、
HubSpot、Asana、Intercom、Linear 等。
它不仅索引这些应用中的内容（文档、消息、工单、代码等），
还理解应用之间的关系和工作流。
例如，Glean 知道一个 JIRA 工单引用了哪些 Confluence 页面，
一个 Slack 讨论涉及了哪些 Google Docs，
一个 GitHub PR 关联了哪些 Linear 任务。
这种跨应用的关系理解是 Glean 区别于简单关键词搜索的关键。

\textbf{知识图谱层}（Knowledge Graph）。
Glean 为每个组织构建一个定制化的知识图谱，
映射组织的三个维度：

\begin{itemize}[noitemsep]
  \item \textbf{内容维度}：文档、消息、工单、代码、视频——
        组织中所有数字化内容的语义索引。
  \item \textbf{人员维度}：谁负责什么项目？谁是某个技术领域的专家？
        谁最近在活跃地使用某个工具？
        组织架构、协作关系和专业知识分布的映射。
  \item \textbf{交互维度}：谁最近在讨论某个话题？
        哪些文档被频繁引用？哪些 Slack 频道在密集交流？
        组织活动模式和知识流动方向的实时追踪。
\end{itemize}

这个知识图谱是 Glean 区别于通用 AI 聊天机器人的核心资产——
它理解的不仅是``语言''，而是``组织''。
ChatGPT 知道``什么是 Kubernetes''，
但 Glean 知道``在你的公司里，谁负责 Kubernetes 集群管理，
上周的集群迁移文档在哪里，
以及你的同事 Alex 昨天在 Slack 里关于这次迁移说了什么''。

\textbf{AI 层}（AI Engine）。
基于连接层的数据和知识图谱的结构化理解，
Glean 提供三种交互模式：

\begin{enumerate}[noitemsep]
  \item \textbf{AI 搜索}：回答``我在哪里可以找到 X？''——
        跨所有企业应用的统一搜索，权限感知，个性化排序。
  \item \textbf{AI 助手}：回答``根据我们的内部知识，X 的答案是什么？''——
        基于企业内部文档生成回答，附带来源引用，确保可追溯性。
  \item \textbf{AI Agent}：执行``帮我完成 X 这个跨系统的工作流''——
        例如``帮我把这个 bug report 创建为 JIRA 工单，
        通知相关的工程师，并在 Slack 频道发布更新''。
\end{enumerate}

\subsubsection{权限感知：安全设计的核心}
Glean 的一个核心设计原则是\textbf{权限感知}（permission-aware）。
每一次 AI 生成的回答都严格遵守源系统的访问控制——
如果某个用户在 Confluence 中无权查看某个页面，
那么 Glean 的 AI 也不会将该页面的内容暴露给他。
这个看似简单的设计，实际上是 Glean 在企业市场制胜的关键。

正如前文所述，Microsoft Copilot 在数据权限方面遇到了挑战——
``幽灵数据暴露''问题迫使企业在部署前进行大规模权限清理。
而 Glean 从第一天起就将权限治理作为架构的核心，
而非事后补丁。
Glean 的平台``设计为确保每一个操作都经过认证，
并受企业安全护栏的约束''——
用 Arvind Jain 自己的话说，
``客户保持对其信息的完全控制权。''

这一安全设计的另一个层面是 Glean 的\textbf{数据主权}承诺。
Glean 不会使用客户数据来训练通用模型——
每个客户的知识图谱和 AI 模型都是独立的，
不存在``你的数据被用来改善竞争对手的搜索体验''的风险。
在后 GDPR 时代，这种数据主权承诺对欧洲和亚太企业客户尤为重要。

\subsubsection{融资与估值：9 个月翻了 57\%}
Glean 的融资历程堪称企业 AI 赛道的缩影：

\begin{itemize}[noitemsep]
  \item 2024 年 9 月：Series E，融资 2.6 亿美元，估值 46 亿美元。
        由 Altimeter 领投，
        投资方包括 Kleiner Perkins、Lightspeed Venture Partners、
        Sequoia Capital、DST Global、General Catalyst、ICONIQ、IVP 等顶级 VC。
  \item 2025 年 6 月（仅 9 个月后）：Series F，
        融资 1.5 亿美元，估值 72 亿美元。
        由 Wellington Management 领投，
        新投资方包括 Khosla Ventures、Bicycle Capital、
        Geodesic Capital 和 Archerman Capital。
        现有投资者几乎全部跟投，包括 Altimeter、Capital One Ventures、
        Citi、Coatue、DST Global、General Catalyst、ICONIQ、IVP、
        Kleiner Perkins、Latitude Capital、Lightspeed Venture Partners、
        Sapphire Ventures 和 Sequoia Capital。
  \item 2026 年 2 月：Glean 确认 Series F 轮（Crunchbase/PitchBook 将此次融资
        归类为 2026 年的确认轮次），估值维持在 72 亿美元。
        至此，Glean 的总融资额累计超过 6 亿美元。
\end{itemize}

9 个月内估值从 46 亿美元跃升至 72 亿美元（增长 57\%），
这在已经降温的 2025 年 VC 市场中极为罕见。
Crunchbase 评论道：``Glean 的估值跃升发生在距上次融资仅九个月之后，
在当前的融资环境中尤为引人注目。''

\subsubsection{增长指标：ARR 破亿的飞速成长}
支撑这一估值的是 Glean 惊人的增长数据：

\begin{itemize}[noitemsep]
  \item \textbf{ARR 突破 1 亿美元}。Glean 在 2025 年初达到了这一里程碑，
        且增速仍在加速。达到 1 亿美元 ARR 的速度
        使其跻身企业软件历史上增长最快的公司之列。
  \item \textbf{年处理超过 1 亿次 Agent 行为}。
        这个数字极为重要——它意味着 Glean 已经从``搜索工具''进化为``执行平台''。
        用户不仅通过 Glean 查找信息，
        还通过 Glean 的 Agent 完成跨系统的操作——
        创建工单、发送通知、生成报告、更新文档等。
  \item \textbf{员工规模约 1000 人}（截至 2026 年初），
        涵盖工程、产品、销售和客户成功团队。
        这一人员规模对应超过 1 亿美元的 ARR，
        意味着每员工收入约 10 万美元——
        对于一家仍处于高速增长阶段的公司来说，
        效率相当可观。
  \item \textbf{国际扩张}。Series F 的资金被明确用于
        ``加速产品开发、扩展合作伙伴生态系统和支持国际增长''，
        表明 Glean 正在从北美市场向全球扩展。
\end{itemize}

\subsubsection{Glean 的战略启示}
Glean 的成功揭示了企业 AI 创业的几个关键洞察：

\textbf{搜索是企业 AI 的最佳切入点}。
企业员工每天花费 20--30\% 的时间在``寻找信息''上。
从搜索切入，Glean 解决了一个高频、高痛点的需求，
并在此基础上自然扩展到 AI 助手和 AI Agent。
这比试图一步到位地推出``企业 AI 平台''要务实得多——
搜索是一个``每个人每天都需要''的功能，
而``AI Agent''对大多数企业员工来说仍是一个抽象概念。
``从搜索到 Agent''的渐进路径，
确保了产品在每个阶段都有清晰的价值主张。

\textbf{知识图谱是真正的壁垒}。
Glean 为每个客户构建的组织知识图谱是一种``结构化的上下文''，
它使 Glean 的 AI 能够提供高度个性化、组织感知的回答。
这是通用 LLM（如直接使用 GPT-4 或 Claude）所无法复制的。
一旦企业部署了 Glean 并积累了知识图谱，
迁移成本就变得极高——你的竞争对手可以复制你的产品功能，
但无法复制你已经为 1000 家企业构建的 1000 个独特的知识图谱。
这是典型的数据网络效应——
每一个新的用户交互都在丰富知识图谱，
使产品对现有用户更有价值。

\textbf{``从 Google 走出来''的创始人优势}。
Arvind Jain 不仅带来了搜索技术的深厚积累，
还带来了对``企业级搜索需要什么''的第一手理解。
Google 内部的知识检索系统是世界上最复杂的企业搜索系统之一——
它需要处理 18 万员工、数十亿文档、数百个内部工具。
Jain 将这种经验外部化、产品化，
创造了一个通用的企业搜索平台。
这种``从巨头内部走出来的创始人''模式
在企业 AI 赛道中反复出现——
ThoughtSpot 的创始人来自 Google，
Moveworks 的创始团队来自 Facebook 和 Google，
Glean 本身就是这一模式的典范。

\subsection{Guru AI 与 Aisera：知识管理的其他路径}

除 Glean 外，企业搜索与知识管理赛道还有几个值得关注的参与者，
它们代表了不同的产品哲学和市场定位。

\textbf{Guru} 的定位是``AI 驱动的企业知识中心''。
与 Glean 的``搜索一切''模式不同，
Guru 更强调``策展式''（curated）的知识管理——
核心功能是将散落在各处的企业知识整合成一个可信赖的、
经过验证的知识库。
Guru 的差异化在于其``知识验证''（Knowledge Verification）机制——
内容由人类专家定期审核和更新，
过期或不准确的知识会被自动标记和下架。
AI 则帮助将这些经过验证的知识
在正确的时间推送给正确的人。

这种``策展式''的方法在某些场景中优于``搜索一切''——
特别是在新员工 onboarding、客户服务知识库、
和合规性要求高的行业（如医疗、金融）中，
``被验证过的正确答案''比``AI 从海量内容中推断的答案''更有价值。
Guru 的核心客户群体是销售团队和客户服务团队，
这些团队需要快速、准确的内部知识支持来响应外部需求。

\textbf{Aisera} 则专注于 AI 驱动的服务体验，
包括 IT 服务台、HR 服务台和客户服务。
值得注意的是，Automation Anywhere 在 2025 年收购了 Aisera，
将其 AI 服务台和企业搜索能力整合进了自己的自动化平台。
这次收购反映了企业 AI 赛道的一个重要趋势：
\textbf{AI 搜索和 AI 自动化正在融合}——
``找到信息''和``基于信息执行行动''之间的界限正在模糊。
用户不想先搜索``如何重置密码''的文档，
然后再手动按照步骤操作；
他们想直接告诉 AI ``帮我重置密码''，
AI 自动执行所有步骤。
这就是``从搜索到行动''的进化方向。

\subsection{Moveworks 的结局：被 ServiceNow 收购的启示}

如前文所述，Moveworks 在 2025 年 3 月被 ServiceNow 以 28.5 亿美元收购。
这家成立于 2016 年的公司，
曾以``用 AI 自动解决企业 IT 问题''为核心定位，
融资超过 3 亿美元，估值在 2021 年达到 21 亿美元。

Moveworks 的被收购验证了企业 AI 搜索/助手赛道的一个重要规律：
\textbf{独立的企业 AI 助手公司面临``平台夹击''风险}。
当 ServiceNow、Salesforce、Microsoft 等平台巨头
都在构建自己的 AI 助手时，
独立的 AI 助手公司只有两条出路：

\begin{enumerate}[noitemsep]
  \item 找到一个平台无法覆盖的差异化定位
        （如 Glean 的\textbf{跨平台}搜索——它不属于任何一个平台生态，
        而是连接所有平台的``桥梁''）。
  \item 被平台收购，成为平台能力的一部分
        （如 Moveworks 被 ServiceNow 收购，
        成为 ServiceNow AI 体验的前端层）。
\end{enumerate}

没有第三条路。
试图在平台巨头的核心领地（如 IT 服务管理）
与平台正面竞争的独立公司，
几乎不可能长期存活。

对 Moveworks 的投资者来说，
这次收购是一个不错的回报——
3 亿美元投资换来 28.5 亿美元收购价，约 9.5 倍回报。
但对 Moveworks 的创始团队来说，
这也标志着一家独立企业 AI 公司愿景的终结。
这种``成功的收购 = 独立性的丧失''的悖论，
是企业 AI 创业者必须面对的战略现实。


%% ============================================================
\section{AI 数据分析}
\label{sec:ent-analytics}
%% ============================================================

数据分析是企业 AI 渗透最深、商业验证最充分的垂直领域之一。
传统的 BI（Business Intelligence）工具——Tableau、Power BI、Looker——
要求用户掌握 SQL 或至少理解数据模型才能生成报表。
即使有了``拖拽式''的可视化界面，
数据分析在大多数企业中仍然是少数``数据分析师''的专利——
广大的``业务用户''（销售经理、市场总监、运营主管）
只能提需求、等报表，无法自主获取洞察。
AI 的介入正在将数据分析的门槛降至自然语言。

\subsection{ThoughtSpot：``Agentic Analytics''的开创者}

ThoughtSpot 成立于 2012 年，
由前 Google 工程师 Amit Prakash 和 Ajeet Singh 创立。
核心理念是``像 Google 搜索一样查询企业数据''——
用户只需用自然语言提问
（如``上个季度华东区最畅销的产品是什么？''或
``过去三个月客户流失率的趋势如何？''），
ThoughtSpot 就能自动生成 SQL 查询、执行分析并返回可视化结果。

截至 2026 年，ThoughtSpot 已融资 6.74 亿美元
（包括 2021 年 11 月的 1 亿美元 Series F），
估值 42 亿美元，员工约 763 人。
其客户包括 T-Mobile、BT（英国电信）、Snowflake、
Daimler（戴姆勒）、CVS、OpenTable、Capital One、
Huel 和 Verisk Analytics 等。

其最新的战略定位是``Agentic Analytics Platform''——
从``帮人查数据''升级为``让 AI Agent 自主发现洞察并触发行动''。
这一定位转变不仅仅是品牌策略，
而是反映了产品能力的根本性进化：
ThoughtSpot 的 AI Agent 现在能够主动监控数据变化、
自动生成异常报告、并建议后续分析方向。

\subsubsection{ThoughtSpot 全球研究：数据洞察的马太效应}
ThoughtSpot 2026 年 3 月发布的``分析新运营模式''报告
（基于全球 1200 名数据和业务领导者的调查）揭示了几个关键趋势：

\begin{itemize}[noitemsep]
  \item 74\% 的企业计划在三年内达到``完全 GenAI 成熟度''。
  \item 但仍有近 40\% 的企业需要等待超过 24 小时才能获得单个数据洞察，
        24\% 的企业甚至需要等待超过一周。
        ThoughtSpot 将此称为``遗留延迟危机''（Legacy Latency Crisis）。
  \item 在``领先型''组织中（已经将 GenAI 和 Agentic AI 用例运营化的企业），
        53\% 的员工能够\textbf{即时}获取可信赖的数据洞察。
        与之形成鲜明对比的是，``落后型''组织中这一比例不足 10\%。
  \item 这种 AI 成熟度差距正在形成``马太效应''——
        AI 成熟的企业能更快地从数据中提取价值，
        获得更高的投资回报，从而投入更多资源于 AI，
        进一步拉大与竞争对手的差距。
        ThoughtSpot 将此描述为``AI 成熟度是一个自我增强的飞轮''。
\end{itemize}

\subsection{Sigma Computing：电子表格与云数据仓库的 AI 桥梁}

Sigma Computing 代表了另一种 AI 数据分析的思路。
它将\textbf{电子表格界面}与\textbf{云数据仓库}
（如 Snowflake、BigQuery、Databricks）直连，
使业务用户能够在熟悉的电子表格操作范式下直接查询 PB 级数据，
无需学习 SQL 或任何新工具。
AI 功能被叠加在电子表格交互之上——
用户可以用自然语言描述想要的分析，
Sigma 自动生成对应的计算逻辑和可视化。

这种``电子表格 + 云仓库 + AI''的三位一体模式，
对 Excel 重度用户（也就是绝大多数企业员工）极具吸引力。
它的核心洞察是：\textbf{不要试图改变用户的习惯，
而是在用户已有的习惯上叠加 AI 的力量}。
全球有超过 10 亿的 Excel 用户，
``让这些用户在电子表格中直接查询云端的 TB 级数据''
是一个巨大的市场机会。

\subsection{数据分析的 Agent 化趋势}

数据分析领域的 Agent 化趋势可以概括为三个阶段：

\begin{enumerate}[noitemsep]
  \item \textbf{2023--2024 年：Text-to-SQL}。
        用自然语言生成 SQL 查询，但用户仍需理解结果并手动操作。
        典型代表：早期的 ThoughtSpot Sage、AI2sql、SQLChat。
  \item \textbf{2025 年：AI 辅助分析}。
        AI 不仅生成查询，还解释结果、建议可视化方式、
        标记异常值、提供分析备选方案。
        典型代表：ThoughtSpot AI、Sigma AI、Mode AI。
  \item \textbf{2026 年及以后：自主分析 Agent}。
        AI Agent 能够独立执行端到端的数据分析任务——
        从提出假设、设计分析方案、执行查询、
        解释结果到生成行动建议，全程无需人类干预。
        这一阶段的标志是 AI 不仅能``回答问题''，
        还能``提出问题''——
        主动发现数据中人类尚未注意到的模式和异常。
\end{enumerate}

这三个阶段的演进再次印证了本章的核心主题：
企业 AI 正在从``辅助''走向``自主''。

\subsection{数据分析 AI 的竞争格局与赛道总结}

AI 数据分析赛道的竞争格局可以分为三层：

\textbf{平台层}。Microsoft Power BI + Copilot 和 Google Looker + Gemini
凭借其云平台的分发优势占据最大的市场份额。
对于已经在使用 Azure 或 GCP 的企业来说，
使用 Power BI 或 Looker 几乎是``免费''的——
它们捆绑在云服务套件中。
但这些工具的 AI 能力相对浅层，
更多是``在现有 BI 产品上加了一个聊天窗口''。

\textbf{挑战者层}。ThoughtSpot、Sigma Computing、Mode 等公司
提供更深度的 AI 分析能力，
且作为独立产品不受制于特定的云平台生态。
ThoughtSpot 的``Agentic Analytics''定位使其在 AI 深度上
领先于平台层的竞争对手。
但这些公司面临的挑战是：
如何在 Power BI 和 Looker 的``免费捆绑''策略下生存？
答案通常是：服务于对数据分析有更高要求的客户——
那些需要``AI 主动发现洞察''而非仅仅``AI 帮我写 SQL''的企业。

\textbf{新兴层}。Databricks 的 AI/BI Dashboard、
Snowflake 的 Cortex Analyst 等数据平台公司
正在从``数据仓库''向``数据分析''上溯。
这对 ThoughtSpot 等中间层公司构成了``两面夹击''的威胁——
上有 Microsoft/Google 的平台压力，
下有 Databricks/Snowflake 从数据层向上渗透。
ThoughtSpot 的应对策略是深化 AI 能力
（从 Text-to-SQL 到 Agentic Analytics），
使自己不可被``附带''的 AI 功能所替代。

值得关注的一个趋势是``嵌入式分析''（Embedded Analytics）。
越来越多的 SaaS 产品将数据分析能力
直接嵌入到自己的应用中——
Salesforce 内置了 Einstein Analytics，
ServiceNow 内置了 Performance Analytics，
HubSpot 内置了 AI 报告工具。
这意味着用户不需要打开一个独立的 BI 工具来分析数据，
数据分析发生在他们日常使用的应用内部。
ThoughtSpot 和 Sigma Computing 都提供了``嵌入式''版本，
允许其他 SaaS 公司将 AI 分析能力嵌入自己的产品——
这是它们应对``SaaS 内置分析''趋势的关键策略。


%% ============================================================
\section{AI 网络安全}
\label{sec:ent-security}
%% ============================================================

网络安全是 AI 在企业领域最紧迫的应用场景之一，
也是``AI vs AI''对抗最为激烈的领域。
攻击者已经在使用 AI 生成钓鱼邮件、自动发现漏洞、
规避检测系统、甚至发动自主性的网络攻击。
防守方如果不以 AI 对 AI，就注定处于劣势。
AI 网络安全市场正在以超过 25\% 的年复合增长率扩张，
成为企业 AI 领域增长最快的细分市场之一。

\subsection{CrowdStrike Charlotte AI：端点安全的智能中枢}

CrowdStrike 是全球领先的端点安全公司，
其 Falcon 平台保护着数百万台企业终端设备——
笔记本电脑、服务器、云工作负载和移动设备。
Charlotte AI 是 Falcon 平台的 AI 核心引擎，
提供三层递进的 AI 能力：

\textbf{检测分诊}（Detection Triage）。
Charlotte AI 能够自动评估每一个安全告警的严重程度和可信度，
将安全分析师从``告警疲劳''中解放出来。
一个中型企业的 SOC（安全运营中心）每天可能收到数万条告警，
其中 95\% 以上是误报或低优先级事件。
CrowdStrike 声称 Charlotte AI 的告警分诊准确率达到了``接近人类专家''的水平，
同时处理速度提升了数十倍——
将每条告警的分诊时间从分钟级缩短到秒级。

\textbf{调查辅助}（Investigation Assist）。
当安全事件需要深入调查时，
Charlotte AI 能够自动关联来自不同数据源的信号——
端点日志、网络流量、身份认证记录、云工作负载数据——
并生成结构化的调查报告。
安全分析师可以用自然语言提问：
``这个可疑进程是否与上周的其他告警相关？
它的行为模式匹配了哪些已知的 MITRE ATT\&CK 技术？''
Charlotte AI 会自动执行跨数据源的关联分析，
并以人类可读的叙事形式呈现结果。

\textbf{Charlotte AI AgentWorks}（2026 年推出）。
这是 CrowdStrike 的最新功能，也是最具颠覆性的一个——
允许安全团队构建、编排和扩展 AI Agent，
实现安全运营的自动化。
Charlotte AI AgentWorks 包括 Charlotte Agentic SOAR
（安全编排、自动化和响应），
代表了安全运营从``人类主导、AI 辅助''
到``AI 主导、人类监督''的范式转变。
在这一范式下，AI Agent 能够自主执行完整的事件响应流程——
检测威胁、评估风险、收集证据、执行遏制措施——
人类分析师的角色从``执行者''转变为``监督者''和``决策者''。

CrowdStrike 已获得 ISO/IEC 42001:2023 AI 治理认证
（全球首批获得该认证的安全厂商之一），
并且 Charlotte AI 通过了 FedRAMP 认证，
使其能够服务于美国联邦政府机构——
这是一个对安全要求极高、准入壁垒极强的高价值市场。

\subsection{SentinelOne Purple AI：自主安全的先行者}

SentinelOne 是 CrowdStrike 在端点安全领域的主要竞争对手，
其差异化定位是``自主安全''（Autonomous Security）。

SentinelOne 的 AI 引擎被称为 Purple AI，
核心能力是将安全运营中心（SOC）中原本需要人类分析师
手动执行的调查和响应流程\textbf{完全自动化}。
Purple AI 的三大能力：

\begin{itemize}[noitemsep]
  \item \textbf{自然语言安全查询}：
        安全分析师可以用自然语言查询安全数据
        （类似于安全领域的 ThoughtSpot），
        无需编写复杂的 SIEM 查询语法。
  \item \textbf{自动关联与解释}：
        AI 自动将看似独立的安全事件关联成完整的攻击链，
        并以``故事线''（storyline）的形式呈现，
        使分析师能够快速理解攻击的全貌。
  \item \textbf{自主响应}：
        根据分析结果自动执行响应动作——
        如隔离受感染设备、封锁可疑 IP、
        撤销被盗用的凭证、回滚被修改的文件。
\end{itemize}

SentinelOne 的``自主''理念在行业中引发了激烈辩论。
支持者认为，面对每天数百万条告警和日益复杂的攻击手段，
人类分析师根本无法跟上速度。
2026 年的 CrowdStrike 全球威胁报告显示，
高级攻击者的``突破时间''（breakout time）——
从初始入侵到横向移动的时间——已缩短到分钟级。
在这种速度下，等待人类分析师审查和批准响应动作
可能意味着攻击已经得逞。

批评者则担忧，让 AI 自主执行安全响应
可能造成严重的误伤——
在网络安全领域，将合法流量误判为攻击并封锁
（false positive + auto-response）
可能导致关键业务中断，
其代价可能远高于实际的安全事件。

\subsection{Vectra AI：网络检测与响应的 Gartner 领导者}

如果说 CrowdStrike 和 SentinelOne 聚焦于\textbf{端点}安全，
Vectra AI 则专注于\textbf{网络}检测与响应
（NDR, Network Detection and Response）——
监控和保护网络层面的通信流量。

Vectra AI 在 2025 年被 Gartner 评为
网络检测与响应魔力象限（Magic Quadrant for NDR）的``领导者''（Leader），
且是 MITRE D3FEND 框架中被引用次数最多（\#1）的安全厂商。
公司拥有 35 项以上的 AI 相关专利。

Vectra AI 的核心技术是 Attack Signal Intelligence——
一种 AI 驱动的攻击信号识别系统，
能够跨越数据中心、多云环境、身份认证系统、SaaS 应用、
边缘设备、IoT/OT 和 AI 基础设施，
持续观察、识别和响应网络攻击行为。

Vectra AI 的独特价值在于其\textbf{行为分析}能力。
传统的安全检测依赖于``签名匹配''（signature-based）——
将网络流量与已知攻击模式的特征库进行比对。
这种方法对已知攻击有效，
但对零日攻击（zero-day）、高级持续性威胁（APT）
和供应链攻击等新型威胁无能为力——
因为这些攻击没有``已知签名''。

Vectra AI 的 AI 模型不是寻找``已知的坏模式''，
而是学习``正常的行为基线''
（网络流量的正常模式、用户的正常访问行为、应用的正常通信模式），
然后标记任何偏离基线的异常行为。
这使其能够检测到传统方法遗漏的隐蔽攻击——
例如一个被攻击者控制的合法账户在凌晨 3 点访问了它从未访问过的数据库。

\subsubsection{AI 网络安全的``军备竞赛''}

AI 在网络安全领域催生了一场独特的``军备竞赛''：

\textbf{攻击侧}——AI 被用于：
生成更逼真的钓鱼邮件（传统的``拼写错误''识别法已经失效——
AI 生成的钓鱼邮件在语法和风格上与真实邮件无异）；
自动发现软件漏洞（AI fuzzing 工具能在数小时内
发现人类安全研究员需要数周才能找到的漏洞）；
生成变种恶意代码以规避特征检测；
甚至模拟合法用户行为以绕过行为分析系统。

\textbf{防守侧}——AI 被用于：
实时分析海量安全日志（处理速度是人类分析师的千倍以上）；
自动关联跨系统的攻击信号（连接端点、网络、身份、云等多个数据源）；
预测攻击者的下一步行动（基于 MITRE ATT\&CK 框架的攻击链推理）；
自主执行响应动作（在毫秒级时间内隔离威胁）。

\textbf{治理层}——AI 安全系统本身也成为攻击目标：
对抗性攻击（adversarial attack）可以欺骗 AI 检测模型；
提示注入（prompt injection）可以操纵基于 LLM 的安全助手；
数据投毒（data poisoning）可以在训练阶段破坏 AI 模型的判断力。
这催生了``AI 安全的安全''（securing the security AI）这一新兴子领域。

这场军备竞赛的结果是：
AI 网络安全不再是``可选项''，而是``必选项''。
任何不具备 AI 能力的安全产品，
都将在这场不对称的攻防博弈中逐渐失去竞争力。

\subsubsection{AI 网络安全的市场格局}

从竞争格局来看，AI 网络安全市场呈现出明显的分层结构：

\textbf{第一梯队：安全巨头的 AI 化}。
CrowdStrike（Charlotte AI）、SentinelOne（Purple AI）、
Palo Alto Networks（Cortex XSIAM）和 Microsoft（Security Copilot）
是这一梯队的代表。它们拥有庞大的客户基础和海量的安全数据，
AI 功能的添加是建立在已有优势之上的``强者愈强''策略。
其中，CrowdStrike 和 SentinelOne 作为``端点安全双雄''，
在 AI 驱动的端点检测与响应（EDR）领域形成了寡头竞争格局。

\textbf{第二梯队：AI 原生安全公司}。
Vectra AI（网络检测）、Abnormal Security（邮件安全）、
Intezer（恶意软件分析）等公司从成立之初就以 AI 为核心。
它们在各自的细分领域（NDR、邮件安全、SOC 自动化）
建立了深厚的 AI 技术壁垒。

\textbf{新兴力量：Agentic SOC 创业公司}。
2025--2026 年涌现了一批专注于``Agentic SOC''的创业公司，
试图用 AI Agent 完全替代 SOC 分析师的调查和响应工作。
这些公司声称能够将 SOC 的平均调查时间从数小时缩短到数分钟，
同时处理远超人类团队能力的告警量。
然而，安全领域的``自主决策''门槛极高——
一次错误的自动响应可能导致关键业务系统中断。
因此，这些公司面临的最大挑战不是技术，而是客户信任。

\textbf{``AI 安全的安全''新兴子领域}。
随着企业部署越来越多的 AI 系统，
保护这些 AI 系统本身免受攻击成为一个新的安全需求。
LLM 安全（防止提示注入、数据泄露、模型投毒）、
AI 供应链安全（确保训练数据和模型来源的可信性）、
AI 治理和合规（确保 AI 系统符合 EU AI Act 等法规要求）
正在形成一个全新的安全子市场。
这个子市场在 2024 年几乎不存在，
到 2026 年已经涌现出数十家创业公司。


%% ============================================================
\section{AI 供应链与物流}
\label{sec:ent-supply-chain}
%% ============================================================

全球供应链在 2020--2022 年的疫情冲击中暴露出极端的脆弱性。
港口拥堵、芯片短缺、运费飙升、库存失衡——
这些痛点催生了一波 AI 供应链创业浪潮。
虽然 VC 对供应链赛道的热情在 2023--2024 年有所降温
（融资额从 2021 年的 147 亿美元高峰骤降至不足 20 亿美元），
但那些找到了真正 Product-Market Fit 的公司
正在从``概念验证''走向``规模化部署''。

\subsection{o9 Solutions：AI 驱动的``数字大脑''}

o9 Solutions 由 i2 Technologies 联合创始人 Sanjiv Sidhu 和 Chakri Gottemukkala
于 2009 年在达拉斯创立。
公司的核心产品被称为``Digital Brain Platform''——
一个 AI 驱动的集成规划和决策平台，
将需求预测、供应规划、采购分析、财务计划和商业决策
统一在一个共享的数据和模型环境中。

o9 Solutions 的增长数据令人印象深刻：

\begin{itemize}[noitemsep]
  \item 2024 年收入约 1.575 亿美元，同比增长约 31\%；
        2025 年收入预计达到约 2 亿美元，
        延续了双位数的增长轨迹。
  \item 2025 年 Q1 新增客户同比增长 60\%，
        客户覆盖汽车、消费品、时装、零售等行业。
  \item 总融资 5.36 亿美元，估值 37 亿美元（独角兽），
        投资方包括 General Atlantic、KKR 和 BeyondNetZero。
  \item 员工约 2600 人（截至 2026 年），服务覆盖 29 个国家。
\end{itemize}

o9 的差异化在于其\textbf{``数字大脑''架构}——
不是为供应链的某个环节（如需求预测或库存优化）提供点状解决方案，
而是构建一个统一的决策平台，
将供应链的所有环节和所有相关方
整合到一个共享的数据和模型环境中。

这种``全链路''的方法论在理论上极具吸引力，
但在实践中面临巨大的实施挑战——
它要求客户重构现有的规划流程和数据架构，
部署周期长、变更管理复杂。
o9 的客户主要是大型跨国企业，
这些企业有能力承担长周期的数字化转型投入。

\subsection{Altana AI：全球贸易的``信任网络''}

Altana AI 的定位与传统供应链软件截然不同。
它不是一个``帮你管好自己的供应链''的工具，
而是一个``帮你理解全球供应链''的网络。

Altana 构建了全球最大的结构化供应链数据集，
连接买方、供应商、物流商和政府机构。
其客户包括 Boston Scientific、Maersk（马士基）、L.L.Bean 等商业企业，
以及美国海关与边境保护局（US Customs and Border Protection）、
英国国防部和英国商业贸易部等政府机构。

融资方面，Altana 在 2024 年 7 月完成了由 US Innovative Technology Fund
领投的 2 亿美元 Series C，总融资约 3.24 亿美元。
截至 2025 年底，年收入约 4560 万美元，员工 214 人。

Altana 的战略意义在于它处于\textbf{贸易合规与国家安全}的交叉点。
在中美脱钩、供应链``友岸外包''（friend-shoring）
和出口管制日益收紧的地缘政治背景下，
``了解你的供应链''不再只是效率问题，而是合规和安全的刚需。
企业需要知道：我的三级供应商中是否有受制裁实体？
我的零部件是否经过了被禁止的国家或地区？
我的供应链是否符合欧盟碳边境调节机制（CBAM）的要求？
这些问题只有 Altana 这样拥有全球供应链图谱的平台才能回答。

\subsection{FourKites AI：实时可见性平台}

FourKites 专注于供应链的\textbf{实时可见性}——
追踪在途货物的位置、预计到达时间和潜在延误。
其 AI 引擎基于历史运输数据和实时信号
（天气、交通、港口拥堵、海关状态等）
预测交货时间，准确率在过去几年持续提升。

FourKites 代表了供应链 AI 的一个重要子类别：
不是``规划''供应链（如 o9 Solutions），
也不是``理解''供应链（如 Altana AI），
而是``观察''供应链——实时追踪物流的物理流动。
在一个充满不确定性的世界中，
仅仅知道``货物在哪里、什么时候到''这一看似简单的信息，
对企业的库存管理、生产计划和客户承诺就有巨大价值。

FourKites 的 AI 预测能力在过去三年中持续提升，
这得益于其不断积累的历史运输数据——
每多跟踪一次运输，AI 对类似路线的预测就更准确。
这是供应链可见性领域的数据飞轮效应。

\subsection{供应链 AI 的赛道总结}

将三家代表性公司放在一起，可以清晰地看到供应链 AI 的三个细分赛道：

\begin{itemize}[noitemsep]
  \item \textbf{规划优化}（o9 Solutions，37 亿美元估值）：
        ``未来应该怎么做？''——
        需求预测、供应规划、库存优化。
        AI 在这里的角色是\textbf{决策者}。
  \item \textbf{可见性与追踪}（FourKites）：
        ``现在发生了什么？''——
        实时物流追踪、到达时间预测、异常预警。
        AI 在这里的角色是\textbf{观察者}。
  \item \textbf{合规与风险}（Altana AI，3.24 亿美元融资）：
        ``我的供应链是否安全合规？''——
        供应商风险评估、制裁筛查、碳排放追踪。
        AI 在这里的角色是\textbf{审计者}。
\end{itemize}

值得注意的是，虽然 VC 对供应链赛道的融资热度
从 2021 年的 147 亿美元骤降至 2024 年的不足 20 亿美元，
但这种``退潮''恰恰是市场成熟的标志——
泡沫期的``概念公司''被淘汰，
留下的是真正找到 PMF 的公司。
o9 Solutions 的 60\% 新客户增长和
Altana 的 2 亿美元 Series C 融资
证明了这些公司已经走过了``概念验证''阶段，
进入了``规模化增长''阶段。

供应链 AI 的另一个重要趋势是\textbf{多模态数据融合}。
传统的供应链软件主要处理结构化数据（订单、库存、运费）。
新一代供应链 AI 正在融合非结构化数据源——
卫星影像（监测港口拥堵和库存水平）、
新闻和社交媒体（早期预警地缘政治风险和自然灾害）、
IoT 传感器数据（冷链温度监控、集装箱振动检测）。
这种多模态融合使供应链 AI 从``反应式''进化为``预测式''——
不是在问题发生后才响应，
而是在问题出现征兆时就预警和行动。


%% ============================================================
\section{AI 制造与工业}
\label{sec:ent-manufacturing}
%% ============================================================

制造业是 AI 应用最早期的``试验田''之一，
也是从``POC 地狱''中走出来最艰难的领域。
制造环境的物理性、多样性和安全要求
使得 AI 的部署远比在数字世界中复杂——
你不能像升级一个 SaaS 应用那样``热更新''一条生产线上的 AI 模型。

\subsection{Landing AI：吴恩达的``以数据为中心''革命}

Landing AI 可能是全球知名度最高的制造 AI 公司，
原因很简单——它的创始人是 Andrew Ng（吴恩达），
Google Brain 的联合创始人、
Coursera 的联合创始人、
百度的前首席科学家、
也是 Stanford 大学的计算机科学教授。

\subsubsection{从学术到产业的使命}
Ng 在 2017 年创立 Landing AI，
初始愿景是``将 AI 从受益于少数大公司的技术，
变成受益于所有人的技术''。
这一愿景源于他的一个核心观察：
AI 的商业价值高度集中在少数科技巨头（Google、Facebook、Amazon 等）手中，
而占全球 GDP 超过 50\% 的传统行业（制造业、农业、医疗等）
几乎未能从 AI 中获益。

问题的根源不在于模型——深度学习框架（TensorFlow、PyTorch）是免费的，
预训练模型唾手可得。
问题在于\textbf{数据}。
大多数工厂的缺陷样本极度稀少
（一条高质量产线一天可能只有几个缺陷品），
传统的``以模型为中心''（model-centric）的方法
在小数据场景下根本不起作用——
你无法用 50 张缺陷图片训练一个可靠的缺陷检测模型。

由此，Ng 提出了\textbf{``以数据为中心的 AI''}
（Data-Centric AI）理念：
与其在有限数据上调优模型架构（更深的网络、更复杂的损失函数），
不如系统性地提升数据的\textbf{质量}和\textbf{一致性}——
确保标注准确、标注标准统一、数据分布合理。
这一理念看似朴素，实则具有深远影响。
它后来发展成了一场运动，
影响了整个 AI 行业对数据工程的重视程度。

\subsubsection{LandingLens：制造质检的民主化}
Landing AI 的旗舰产品 LandingLens 是一个视觉检测平台，
将 AI 和深度学习应用于制造质量控制。
其核心设计原则是\textbf{无代码/低代码}——
工厂工程师无需编程背景，
只需通过鼠标点击就能在数小时内构建高级 AI 视觉检测模型，
而非传统方法所需的数周或数月。

LandingLens 的工作流程是：
工厂工程师拍摄少量产品图片（正常品和缺陷品），
在 LandingLens 平台上进行标注（用鼠标框出缺陷位置），
平台自动利用预训练模型和数据增强技术训练一个定制化的检测模型，
模型部署到产线边缘设备上实时检测产品。
整个过程不需要编写一行代码。

客户包括 Stanley Black \& Decker（工业工具制造商）
和 Foxconn（富士康，全球最大的电子产品代工制造商）等。

\subsubsection{融资与战略合作}
Landing AI 在 2021 年完成了 5700 万美元的 Series A 融资，
由 McRock Capital 领投，
Intel Capital 和 Samsung Catalyst Fund 参投，
Insight Partners 和加拿大养老基金 CPPIB 也参与了投资。
此后，公司获得了 Pure Storage 和 Snowflake Ventures 的战略投资
（2024 年），以及 ABB Robotics Ventures 的投资（2025 年 9 月）。

ABB 的投资尤为值得关注。
ABB 是全球最大的工业机器人制造商之一，
其投资 Landing AI 的目标是将视觉 AI 与机器人自动化深度融合——
让机器人不仅能``做''（执行动作），
还能``看''（理解环境）和``判断''（做出决策）。
ABB Robotics 总裁 Sami Atiya 表示：
``AI 在机器人领域的需求由更大的灵活性、
更快的部署周期和专业技能短缺所驱动。
与 Landing AI 的合作意味着安装和部署时间
从数周缩短到数小时，
允许更多企业更智能、更快速、更高效地实现自动化。''

\subsubsection{Landing AI 的行业意义}
Landing AI 的故事超越了一家公司的成败。
它代表了一种\textbf{制造 AI 的方法论变革}——
从``请 AI 专家来工厂部署复杂模型''
到``让工厂工程师自己用无代码工具构建 AI''。
这种``民主化''的路径是制造 AI 规模化的关键——
全球有数十万家工厂，每家工厂的产线、产品和缺陷类型各不相同，
不可能为每家工厂配备 AI 专家团队。
只有当工厂自己能够构建和维护 AI 模型时，
制造 AI 才能真正实现``千行百业''的渗透。

\subsection{Instrumental：电子产品制造的 AI 质检}

Instrumental 聚焦于电子产品制造的 AI 质量检测。
其平台利用高分辨率相机和 AI 算法，
在生产线上实时识别 PCB（印刷电路板）、
组件组装、焊接等环节的缺陷。

Instrumental 的差异化在于其\textbf{根因分析}（Root Cause Analysis）能力——
不仅能发现缺陷，还能追溯缺陷的根本原因
（如供应商物料批次差异、设备参数漂移、
操作员行为变化等），
从而帮助工厂从``事后检测''转向``事前预防''。
这种``从检测到预防''的升级，
是制造 AI 从``成本中心''转变为``价值创造''的关键——
发现一个缺陷只是避免了一个坏品的成本，
而预防一类缺陷则可以节省整个批次的返工和报废成本。

\subsection{Sight Machine：工厂数字孪生}

Sight Machine 的核心概念是\textbf{制造数字孪生}
（Manufacturing Digital Twin）——
为工厂的每台设备、每条产线、每个生产流程
创建一个数字化的虚拟副本。
AI 在数字孪生之上运行，
持续分析设备性能、预测故障、优化工艺参数。

Sight Machine 的客户包括多家全球 500 强制造企业。
其价值主张可以用一句话概括：
``你无法优化你看不见的东西。
数字孪生让你第一次能够`看见'整个工厂的运行状态，
AI 则告诉你如何让它运行得更好。''

制造数字孪生的更深层价值在于``跨工厂学习''——
如果一家企业拥有 20 家工厂，
每家工厂的数字孪生数据可以被汇总分析，
识别出``为什么 A 工厂的良品率比 B 工厂高 5\%''的根本原因，
并将最佳实践从一家工厂复制到其他工厂。
这种``组织级学习''是单一工厂优化所无法实现的。

\subsection{制造 AI 的赛道总结}

将本节讨论的三家公司与更广泛的制造 AI 生态系统结合来看，
可以识别出几个关键趋势：

\textbf{从``事后检测''到``事前预防''再到``实时优化''}。
制造 AI 的价值链正在从左向右延伸——
第一阶段是事后检测（在成品中发现缺陷），
第二阶段是事前预防（预测缺陷发生的根因并提前干预），
第三阶段是实时优化（AI 持续调整工艺参数以最大化良品率和产出）。
Landing AI 主要在第一阶段，Instrumental 跨越第一和第二阶段，
Sight Machine 覆盖了全部三个阶段。

\textbf{``边缘 AI''的关键角色}。
制造 AI 不能完全依赖云端——
生产线的节拍时间以毫秒计，
等待云端推理返回结果的延迟是不可接受的。
因此，制造 AI 需要在``边缘''（即工厂本地）部署推理能力。
这对 AI 模型的大小和推理效率提出了严格要求——
模型必须足够小以在工业级边缘设备上运行，
同时足够准确以满足生产质量标准。
NVIDIA 的 Jetson 系列和 Intel 的 OpenVINO
是制造边缘 AI 领域最常用的硬件和软件平台。

\textbf{``数据为中心''方法论的胜利}。
Andrew Ng 提出的``以数据为中心的 AI''理念
已经从 Landing AI 扩散到整个制造 AI 行业。
越来越多的制造 AI 公司意识到：
在数据稀缺的工业场景中，
投入更多精力在数据质量上（标注一致性、数据平衡性、增强策略），
比投入在模型架构上更有效。
这一理念的影响远超制造领域——
它正在改变整个 AI 行业对``数据工程''vs``模型工程''
重要性排序的认知。


%% ============================================================
\section{AI 建筑与房地产}
\label{sec:ent-construction}
%% ============================================================

建筑行业是全球最大的行业之一（全球年产值超过 13 万亿美元），
却也是数字化程度最低的行业之一。
项目超期、超预算是常态而非例外——
麦肯锡的研究显示，大型建筑项目平均超期 20\%、超预算 80\%。
AI 正在从多个角度渗透这个顽固的传统行业——
而当前的数据中心和半导体工厂建设热潮
为建筑 AI 提供了前所未有的市场需求。

\subsection{OpenSpace：360 度影像的 AI 分析}

OpenSpace（成立于 2017 年，总部旧金山）
开发了一个``可视化智能平台''（Visual Intelligence Platform），
核心工作流程是：

\begin{enumerate}[noitemsep]
  \item 工人使用智能手机、360 度相机或无人机在施工现场采集影像；
  \item 影像自动上传至 OpenSpace 云平台，
        与建筑的 BIM（建筑信息模型）数据对齐；
  \item AI 引擎分析影像，自动生成施工进度报告、
        识别与设计图纸的偏差、检测安全隐患、
        量化工程量完成度。
\end{enumerate}

OpenSpace 已融资约 2 亿美元（至 Series D），
年收入约 1500 万美元，员工约 219 人（同比增长 15\%）。
值得注意的是，OpenSpace 收购了英国公司 Disperse.io，
进一步增强了其 AI 影像分析能力。
其平台已与 Procore 和 Autodesk——
建筑行业两个最主要的项目管理和 BIM 平台——实现深度集成。

OpenSpace 的产品设计体现了建筑 AI 的一个核心原则：
\textbf{数据采集的成本必须足够低}。
建筑工人不是技术专家，
让他们使用复杂的扫描设备是不现实的。
OpenSpace 的方案是让工人``戴着 360 度相机正常走路''，
系统自动完成所有的影像处理和分析。
数据采集成本几乎为零，
这是产品能够规模化推广的前提。

\subsection{Buildots：建筑项目的 AI 大脑}

Buildots（成立于 2018 年，总部以色列）
是建筑 AI 领域增长最快的公司之一。
其利用 AI 和计算机视觉技术，
为建筑项目引入了\textbf{绩效驱动的管理方法}——
用客观的 AI 分析替代主观的人工判断，
使项目团队和管理层能够基于准确、无偏见的数据做出决策。

Buildots 声称其技术能够帮助建筑项目\textbf{将延误减少最多 50\%}，
并释放数百万美元的成本节约——
考虑到大型建筑项目动辄数亿美元的预算，
即使只有 1\% 的成本改善也是巨大的价值。

Buildots 在 2025 年 5 月完成了 4500 万美元 Series D 融资，
由 Qumra Capital 领投，
OG Venture Partners、TLV Partners、Poalim Equity、
Future Energy Ventures 和 Viola Growth 参投。
总融资额达到 1.66 亿美元。
公司声称收入连续多年实现三位数增长，
且已签署多个七位数的企业级协议——
这标志着从``项目级别的技术试用''到``企业级别的长期合同''的重要转变。
员工规模约 294 人（同比增长 35\%），
服务覆盖 9 个国家。
Buildots 还收购了 Genda, Inc.，进一步扩展了其技术能力。

CEO Roy Danon 表示：
``实现运营卓越需要时间，需要行业的重大转变，
但潜在的影响是巨大的。
我们已经看到了从按项目采用技术
到签署长期企业协议的转变。
随着企业从被动管理转向主动管理，
行业将释放数十亿美元的成本节约和运营效率。
这就是我们正在建设的未来。''

\subsection{Procore AI 与 Zillow AI}

\textbf{Procore} 是建筑行业最大的项目管理平台之一（公开上市公司），
正在将 AI 嵌入其核心产品。
Procore 的 AI 功能聚焦于\textbf{风险预测}——
基于历史项目数据（数十万个建筑项目的结构化数据集），
预测哪些任务最可能延误、
哪些分包商最可能出现质量问题、
哪些安全事件风险最高。
Procore 的核心优势在于其庞大的数据资产——
作为建筑行业最大的平台，
它拥有其他任何公司都难以复制的行业数据飞轮。

\textbf{Zillow AI} 则代表了 AI 在\textbf{房地产交易}端的渗透。
Zillow 的 Zestimate（AI 房产估值）算法已经成为美国房地产市场的
事实标准参考——即使存在争议，
买卖双方都会在谈判中引用 Zestimate 的估值。
近年来，Zillow 进一步将 AI 应用于：
个性化房源推荐、虚拟看房体验增强、
抵押贷款预审批自动化、以及 AI 驱动的房产市场预测。

Zillow 的案例说明了一个有趣的现象：
AI 不仅能改变建筑的``建造''过程，
还能改变建筑的``交易''过程。
建筑从设计、施工到销售和管理的整个生命周期，
都在被 AI 重新塑造。

\subsection{建筑与房地产 AI 的赛道特征}

建筑和房地产 AI 赛道有几个值得关注的独特特征：

\textbf{极度碎片化的市场}。
全球建筑行业由数百万家中小企业构成——
绝大多数是员工不足 50 人的小型承包商。
AI 产品的销售和部署在这样碎片化的市场中极具挑战。
OpenSpace 和 Buildots 的策略是瞄准大型总承包商和开发商
（这些企业有预算和技术能力采用 AI），
再通过供应链关系向下渗透到分包商。

\textbf{数据中心建设热潮的推动}。
2024--2026 年全球数据中心建设投资超过 2000 亿美元，
由 AI 训练和推理对算力的爆炸性需求驱动。
数据中心项目对工期和质量有极其严格的要求——
每天的延误可能意味着数百万美元的收入损失
（因为 GPU 集群无法按时上线）。
这种``高价值、零容忍延误''的项目特征，
使得建筑 AI 的 ROI 论证变得极其清晰，
成为建筑 AI 采用加速的最大推动力之一。

\textbf{``竣工后''AI 的新兴机会}。
目前大部分建筑 AI 聚焦于``施工阶段''——
进度监控、质量检测、安全管理。
但一个更大的市场机会正在``竣工后''阶段浮现——
建筑运维（Facility Management）的 AI 化。
一栋建筑的施工周期可能只有 2--3 年，
但运维周期长达 30--50 年。
AI 驱动的建筑运维——
智能能耗管理、预测性维护、空间利用率优化——
可能是比施工 AI 更大的长期市场。

\textbf{BIM 数据的战略价值}。
BIM（建筑信息模型）是建筑行业的``数字基因组''——
包含建筑的三维几何模型、材料清单、结构计算、
MEP（机电暖通）系统布局等全部设计信息。
AI 与 BIM 数据的结合打开了巨大的可能性空间——
AI 可以自动检查 BIM 模型中的冲突和错误、
基于 BIM 自动生成施工方案、
将实际施工进度与 BIM 设计进行自动比对。
Autodesk（BIM 领域的领导者）正在积极投资 AI 能力，
OpenSpace 和 Buildots 也都与 Autodesk 进行了深度集成。
控制 BIM 数据的公司，将在建筑 AI 时代拥有核心数据壁垒。


%% ============================================================
\section{AI 农业与能源}
\label{sec:ent-agriculture}
%% ============================================================

如果说网络安全和金融是``数字原生''的 AI 应用场景，
那么农业和能源则代表了 AI 向\textbf{物理世界}最深处的渗透。
在这些领域，AI 需要与传感器、卫星、无人机、
物联网设备和工业控制系统紧密结合，
挑战远超``软件吃软件''的范畴。

\subsection{Arable：田间智能决策}

Arable 开发了一种名为``Mark''的多参数田间传感器，
能够同时测量 40 多种环境和作物指标
（降水量、土壤水分、叶面湿度、光照强度、温度、
蒸散量、光合有效辐射等）。
AI 引擎基于这些实时数据，
为农民提供灌溉建议、病虫害预警、
收获时间预测和产量估算。

Arable 的独特之处在于其``从传感器到决策''的全链路整合。
大多数农业科技公司要么只做硬件（传感器），
要么只做软件（数据分析平台），
Arable 则将两者垂直整合，确保数据质量和分析模型的紧密耦合。
这种垂直整合在早期阶段增加了公司的运营复杂度
（需要同时管理硬件供应链和软件开发），
但在产品成熟后形成了强大的竞争壁垒——
因为数据质量直接决定 AI 模型的准确性，
而只有自己控制传感器才能保证数据质量。

\subsection{Taranis：卫星与无人机的``作物 AI''}

Taranis 利用卫星影像、无人机航拍和人工智能，
为大规模农业提供\textbf{亚厘米级精度}的作物健康监测。
其 AI 能够在数百万张高分辨率影像中
自动识别病害、虫害、杂草、营养缺乏等问题，
精确到单株作物的级别。

Taranis 的价值主张是``精准农业的眼睛''——
在传统农业中，农民依靠肉眼巡视和经验判断来管理作物健康，
这在千亩乃至万亩规模的农场中既不现实也不准确。
Taranis 的 AI 使大规模农场的``精准管理''成为可能——
不是对整块农田统一施药，
而是精确到``第 7 排第 23 株作物有早期叶锈病，
建议局部施用 X 农药 Y 毫升''。
这种精准度不仅提高了防治效果，
还大幅减少了农药使用量——
对环境可持续性和食品安全都有深远意义。

\subsection{SparkCognition/Avathon：工业 AI 平台}

SparkCognition（已于 2024 年更名为 Avathon）
是一家专注于工业 AI 的公司，
由 Amir Husain 于 2013 年在奥斯汀创立，
后将总部迁至硅谷。
前 BP 首席执行官 Lord John Browne 担任董事会主席，
前思科 CEO John Chambers 担任顾问。

Avathon 的核心能力是\textbf{预测性维护}和\textbf{资产优化}——
利用 AI 分析工业设备（涡轮机、压缩机、泵、管道等）的传感器数据，
预测设备故障、优化运行参数、延长设备寿命。
其客户涵盖能源（石油、天然气、可再生能源）、
航空、采矿和制造等行业。
值得一提的是，Avathon 声称驱动了印度超过 20\% 的加油站，
显示了其在能源下游领域的深厚渗透。

Avathon 的融资总额约为 3 亿美元（至 2022 年的 Series D，
融资 1.23 亿美元），
估值约 14 亿美元（独角兽），
投资方包括 Temasek（淡马锡）、March Capital 和 Boeing。

2024 年的品牌重塑标志着公司从
``点状的预测性维护工具''
向``系统级的工业 AI 平台''的战略转型。
新的 Avathon 平台采用了\textbf{系统级方法}（systems-level approach），
不仅监控单个设备，
还理解设备之间的相互依赖关系，
从而实现整个工业系统的全局优化。
例如，在一个石化工厂中，
Avathon 不只是预测某台压缩机何时会故障，
还分析压缩机故障如何影响下游的反应器、分馏塔和管道系统，
从而制定最优的维护和调度策略。

2026 年初，Avathon 与采矿化学品公司 Draslovka 建立了战略联盟，
将 AI 与过程智能结合应用于采矿行业——
这是工业 AI 向更深层、更专业领域渗透的又一例证。

\subsection{Uptake：工业 AI 先行者的经验教训}

Uptake 是工业 AI 领域的先行者之一，
由 Groupon 联合创始人 Brad Keywell 于 2014 年在芝加哥创立。
Uptake 的核心产品是一个工业 AI 平台，
聚焦于重型装备（建筑机械、采矿设备、铁路车辆等）的
故障预测和运维优化。

Uptake 的发展历程本身就是工业 AI 创业的一个教科书式案例——
一个关于``过高期望、过快扩张和艰难调整''的故事。
公司在成立仅两年半时便达到独角兽估值（超过 20 亿美元），
获得了 Caterpillar（卡特彼勒）等工业巨头的投资和背书。
但随后经历了增长放缓、裁员和估值下调。
这反映了工业 AI 创业的一个核心挑战：
\textbf{在重工业领域，销售周期长（6--18 个月）、
部署复杂度高（需要与 OT 系统集成、现场部署边缘设备）、
POC 到生产部署的转化率低（行业平均不到 30\%）。}

Uptake 的经验教训对所有工业 AI 创业者都有警示价值：
工业 AI 的商业成功不仅取决于技术能力，
更取决于对客户组织变革的深度理解。
在重工业领域，最大的阻力往往不是``AI 不够好''，
而是``现场操作人员不信任 AI''或``现有 IT/OT 基础设施无法支撑 AI''。
一个 AI 模型预测``涡轮机将在 72 小时内故障''，
但操作人员说``这台机器我开了 20 年，我比任何 AI 都了解它''——
这种人机信任问题是工业 AI 的``最后一公里''。

\subsection{AI 在能源转型中的角色}

在更宏观的层面上，AI 在能源领域的应用正在从
``优化化石能源生产''转向``加速能源转型''。

\textbf{可再生能源预测}。风力和太阳能发电的间歇性是其最大的挑战。
AI 能够基于天气预报、历史发电数据和电网状态，
更准确地预测可再生能源的发电量，
使电网调度更加高效。
Google DeepMind 的研究显示，
AI 可以将风电预测误差降低 20\% 以上——
这看似微小的改善，对应着数十亿美元的经济价值。

\textbf{电网优化}。
随着分布式能源（屋顶太阳能、家庭储能、电动汽车）的普及，
电网管理变得前所未有地复杂——
从``少数大电厂供电给多数用户''的单向模式，
变成``无数小型产消者互相交换电力''的双向复杂网络。
AI 驱动的电网优化系统能够实时平衡供需、
管理电压波动、优化储能调度。

\textbf{碳排放监测}。
AI 卫星影像分析技术
能够从太空监测甲烷泄漏、碳排放源和森林覆盖变化，
为碳交易和气候合规提供数据支撑。
这一应用领域在欧盟碳边境调节机制（CBAM）
和全球 ESG 披露要求日益收紧的背景下，
正在快速成长。

\subsection{农业与能源 AI 的共性特征}

农业和能源领域的 AI 应用共享几个独特的特征，
使其与``纯数字''的企业 AI 赛道显著不同：

\textbf{物理世界的不可控性}。
在 SaaS 环境中，输入数据是确定的（用户点击、API 返回值）。
在农业和能源环境中，输入数据充满噪声和不确定性——
传感器可能受到极端天气影响而产生错误读数，
卫星影像可能被云层遮挡，
设备振动数据可能因环境变化而产生虚假信号。
AI 模型必须具备更强的鲁棒性和异常处理能力。

\textbf{部署环境的极端性}。
AI 硬件（边缘计算设备、传感器）需要在极端环境中运行——
农田中的高温高湿、油井中的高压高温、
风力涡轮机顶部的强风和盐雾。
这对硬件可靠性和维护成本提出了严苛要求。
一个在数据中心完美运行的 AI 模型，
部署到田间或工厂后可能因为硬件故障而完全失效。

\textbf{ROI 的直接性和可量化性}。
与企业搜索或办公助手的``模糊 ROI''不同，
农业和能源 AI 的价值可以直接量化——
``AI 预测的灌溉方案比传统方案节省了 30\% 的水量和 15\% 的肥料''，
``AI 预测的设备维护避免了一次计划外停机，
节省了 50 万美元的损失和维修费用''。
这种清晰的 ROI 计算是这些领域 AI 采用的最大推动力。

\textbf{行业周期的影响}。
农业和能源都是强周期行业——
粮价、油价、电价的波动直接影响客户的 IT 投入意愿。
在高粮价/高油价时期，客户有更多预算用于 AI 投资；
在低价时期，AI 预算可能是第一个被削减的项目。
这种周期性使得这些领域的 AI 公司需要更强的
现金流管理能力和更灵活的商业模式。

\textbf{监管与合规的独特性}。
农业 AI 涉及食品安全法规、农药使用限制、有机认证标准；
能源 AI 涉及电网安全标准、排放法规、能源交易监管。
这些行业特定的监管要求构成了
通用 AI 平台无法轻易跨越的壁垒——
也是垂直 AI 专家的天然护城河。


%% ============================================================
\section{RPA 遇上 AI}
\label{sec:ent-rpa}
%% ============================================================

RPA（机器人流程自动化, Robotic Process Automation）
是 AI 渗透企业最古老也最务实的路径之一。
传统 RPA 通过模拟人类在屏幕上的点击和输入操作，
将重复性的规则化任务自动化——
登录系统、复制数据、填写表单、生成报告。
它不需要改变底层系统，不需要 API 集成，
只需``看着屏幕操作''就行——
这使得 RPA 成为最容易部署、最快见效的自动化技术。

然而，传统 RPA 有一个根本性的局限：
\textbf{它只能处理结构化的、确定性的任务}。
每一个自动化流程都是一个预定义的脚本，
严格按照``如果 A 则 B、如果 C 则 D''的逻辑执行。
一旦流程中出现非预期的变化
（如 UI 界面更新了一个按钮位置、数据格式出现异常、
弹出了一个意外的对话框），
RPA 机器人就会``宕机''——它不理解``变化''，只会``执行''。

生成式 AI 的崛起正在从根本上改变 RPA 的能力边界。
AI + RPA 的融合创造了一个新范式：
\textbf{Agentic Process Automation}——
不再是机械地执行预定义脚本，
而是让 AI Agent 理解任务目标、
动态适应环境变化、自主做出决策。

2024 年全球 RPA 市场营收达到 38 亿美元，
同比增长 18\%，远超同期基础架构软件 10.7\% 的行业增速。
在 AI 的加持下，RPA 不是在消亡，而是在进化。

\subsection{UiPath：从 RPA 到 Agentic Automation 的转型}

UiPath 是全球最大的 RPA 公司，
由 Daniel Dines 于 2005 年在罗马尼亚创立，
2021 年在纽约证券交易所上市（市值峰值约 350 亿美元），
之后经历了股价大幅回调。
但 FY2026 的财务表现展示了一个成功的``AI 转型''故事。

\subsubsection{FY2026 财务表现}
截至 FY2026（2026 年 1 月 31 日）：

\begin{itemize}[noitemsep]
  \item 年化经常性收入（ARR）达到 18.53 亿美元，同比增长 11\%；
  \item Q4 收入为 4.81 亿美元，同比增长 14\%；
  \item 全年收入约 15.5 亿美元；
  \item GAAP 运营利润为 8000 万美元
        （这对 UiPath 来说是一个重要的里程碑——
        证明其可以在增长的同时实现盈利）；
  \item 非 GAAP 运营利润为 1.5 亿美元；
  \item 非 GAAP 调整后自由现金流为 1.82 亿美元；
  \item 现金及等价物 16.9 亿美元。
  \item 宣布新的 5 亿美元股票回购授权，
        此前的 10 亿美元回购计划已完成。
\end{itemize}

\subsubsection{战略定位的转变}
更重要的是战略定位的转变。
Daniel Dines 在财报电话会上表示：
``随着企业 AI 采用从实验走向规模化部署，
客户越来越需要一个能够以可靠性、治理和规模化方式
执行复杂流程的平台。
通过将确定性自动化、Agentic AI 和企业级编排
整合在一个平台上，
UiPath 提供了企业信赖的、
能够运行关键业务流程的执行层。''

这段话揭示了 RPA 行业的一个深刻转变：
\textbf{RPA 不再只是``自动化''，
而是正在成为 AI Agent 的``执行层''}（execution layer）。
在 Agentic AI 的架构中，
AI Agent 负责``理解''（理解用户意图、理解任务目标）
和``决策''（判断应该做什么、按什么顺序做），
而 RPA 机器人负责``执行''——
在企业系统中点击按钮、填写表单、提交审批、提取数据。
这种``Agent 大脑 + RPA 双手''的分工，
使得 RPA 在 AI 时代不但不会被淘汰，
反而获得了新的战略价值。

UiPath 2026 年趋势报告的核心发现：
78\% 的高管认为他们需要重塑运营模式
才能完全释放 Agentic 自动化的价值。
``单一 Agent''模式已经过时，
``多 Agent 系统''正在成为主流——
多个 AI Agent 协作完成复杂任务，
RPA 提供底层的执行和编排能力，
而``治理即代码''（governance-as-code）
是保持 Agent 对齐、安全和合规的新必备。

UiPath 目前的市值约 70 亿美元（截至 2026 年 3 月），
较 2021 年峰值下跌约 80\%。
市场对 UiPath 的核心疑问是：
在一个 AI Agent 可以直接操控 API 的世界中，
基于``模拟屏幕操作''的 RPA 还有多大的长期价值？
UiPath 的回答是两方面的：
第一，大量遗留系统（legacy systems）没有 API，
RPA 是唯一的自动化手段——
而企业世界中遗留系统的占比远高于人们的想象；
第二，即使有 API，RPA 提供的治理、审计和合规能力
也是直接 API 调用所不具备的——
每一步操作都被记录、可追溯、可审计，
这在受监管行业（金融、医疗、政府）中是硬性要求。

\subsection{Automation Anywhere：Agentic Process Automation 的领导者}

Automation Anywhere 是 UiPath 的主要竞争对手，
成立于 2003 年，总部位于硅谷圣何塞，
由 Mihir Shukla、Ankur Kothari 和 Neeti Mehta Shukla 联合创立。

截至 2026 年，公司年收入约 4.87 亿美元，
员工约 1527 人（同比增长 24\%），
总融资约 13 亿美元（一说 8.4 亿美元，取决于是否计入二级市场交易），
估值约 73 亿美元。
公司在全球 62 个国家运营，服务覆盖面极广。

Automation Anywhere 在品牌定位上比 UiPath 更激进地拥抱 AI。
它自称为``Agentic Process Automation（APA）的领导者''，
并推出了行业首个``流程推理引擎''（Process Reasoning Engine, PRE）——
一种能够让 AI 理解、推理和适应复杂业务流程的技术，
以及四款新的 Agentic AI 解决方案。

值得注意的是，Automation Anywhere 收购了 Aisera——
前文提到的企业 AI 服务台和搜索公司。
这次收购使 Automation Anywhere 获得了 AI 对话界面和企业搜索能力，
再次印证了``AI 搜索 + AI 自动化''融合的趋势。

2025 财年 Q1 报告显示，APA 的预订量增长强劲，
在现有客户基础中的附着率达到 51\%
（即超过一半的现有 RPA 客户已开始采用 APA 功能），
新客户和增购预订实现两位数增长。
EBITDA 超出指引，显示出强健的盈利纪律。

\subsection{来也科技：中国 RPA + AI 的领军者}

在中国市场，来也科技（Laiye）是 RPA + AI 赛道的代表性企业。

来也科技由汪冠春于 2015 年创立，总部位于北京。
公司是\textbf{中国唯一连续五年入选 Gartner RPA 魔力象限}的厂商——
这一成就在中国企业软件领域极为罕见。
截至 2025 年，来也科技在 Gartner 魔力象限中
实现了``愿景完整性''（Completeness of Vision）和
``执行力''（Ability to Execute）双维度的大幅提升，
标志着中国 RPA 技术的全球竞争力持续进阶。

\subsubsection{产品矩阵}
来也科技的产品矩阵包括三大核心板块：

\textbf{数字员工平台}（Digital Worker Platform）。
一个开放的企业级智能体（Agent）平台，
整合了需求管理中心、共享中心、人机协同模块、
开放平台（支持 MCP 协议）和信创安全模块。

\textbf{智能体流程自动化}（Agentic Process Automation）。
包括：流程创造者（可视化流程设计器）、
流程机器人（RPA 执行引擎）、
机器人指挥官（编排和调度中心）、
流程记录者（自动流程发现和文档生成）。

\textbf{智能体文档处理}（Agentic Document Processing）。
基于 AI 的文档理解能力，
包括通用识别、票据识别、文档抽取、
文档分类、文档比对和文档审核。

来也科技的行业解决方案覆盖了中国经济的多个核心领域：
电力、运营商、银行、保险、基金、制造、零售、物流、医疗、
地产、智慧法院、烟草、行政审批、人力社保等。

\subsubsection{AI Agent 产品创新}
来也科技在 2024 年推出了基于大模型的三款 AI Agent 数字员工产品：

\textbf{数字员工开发助手}——核心在于智能化的代码生成机制。
它能够自动解析复杂数据集，
生成与 RPA 产品无缝对接的 Python 代码，
并具备``自愈''功能——自动适应 UI 的变化，
减少对人工维护的依赖。

\textbf{知识管理和问答助手}——通过自动化手段整合分散的文档资源，
经过智能预处理，将文档内容转化为可查询的知识库，
提供 RAG（检索增强生成）驱动的问答能力。

\textbf{文档审核助手}——结合 AI 理解能力和合规规则引擎，
自动审核合同、标书等复杂文档中的关键信息、条款和潜在风险。

\subsubsection{中国 RPA 市场的特殊性}

中国 RPA 市场与全球市场相比有几个显著特征：

\textbf{信创需求}。中国政府的``信息技术应用创新''（信创）战略
要求关键行业（金融、电力、电信、政府）使用国产软件。
这为来也科技等国产 RPA 厂商提供了强大的政策红利——
UiPath 和 Automation Anywhere 在这些受管制行业中
面临严格的准入限制。
来也科技的产品已通过信创适配认证，
能够在国产操作系统（如统信、麒麟）和
国产数据库环境中运行。

\textbf{大模型融合加速}。中国的 LLM 生态（通义千问、文心一言、
讯飞星火等）正在快速与 RPA 融合。
来也科技的平台已支持与多个国产大模型对接，
且已接入 MCP（Model Context Protocol）协议——
这使其数字员工能够与大模型生态更紧密地集成。
大模型通过 MCP 协议调用来也科技的 RPA 能力，
实现``大脑（LLM）+ 手脚（RPA）''的完整闭环。

\textbf{政企市场驱动}。
与美国市场以商业企业为主不同，
中国 RPA 市场有很大一部分来自政府和国有企业。
行政审批自动化、人力社保流程自动化、
法院文书处理自动化等场景，
构成了来也科技收入的重要组成部分。

来也科技累计融资约 2.11 亿美元，
投资方包括 Lightspeed China Partners（光速中国）、
Sequoia Capital China（红杉中国）和 Wu Capital（吾铭资本）等。

来也科技与毕马威中国建立了战略合作伙伴关系
（2021 年签约），共同推动 RPA + AI 在金融行业的落地。
这种``RPA 厂商 + 咨询公司''的合作模式
在中国企业 AI 市场中非常典型——
中国大型企业（特别是金融和国有企业）
习惯于通过咨询公司来采购和实施 IT 解决方案，
直接销售往往效率低下。
来也科技通过与毕马威、德勤等咨询公司的合作，
有效地触达了这些高价值客户。

\subsubsection{RPA 三巨头对比}

将 UiPath、Automation Anywhere 和来也科技放在一起对比，
可以清晰地看到 RPA 行业的全球竞争格局：

\begin{itemize}[noitemsep]
  \item \textbf{UiPath}（ARR 18.5 亿美元，上市公司，市值约 70 亿美元）：
        全球最大，技术深度和产品成熟度最高，
        但增速放缓（11\% ARR 增长），面临估值压力。
        核心市场：北美和欧洲大型企业。
  \item \textbf{Automation Anywhere}（收入约 4.87 亿美元，估值 73 亿美元）：
        品牌定位最激进（``APA 领导者''），
        AI 创新速度快（PRE 引擎、收购 Aisera），
        但仍未盈利，IPO 时间表不确定。
        核心市场：全球大型企业，在中东和亚太有较强存在。
  \item \textbf{来也科技}（融资 2.11 亿美元）：
        中国市场的绝对领导者（唯一连续 5 年入选 Gartner MQ），
        深度绑定信创生态和政企市场，
        但全球化程度有限，收入规模与前两者有数量级差距。
        核心市场：中国政企和金融行业。
\end{itemize}

三者面临的共同挑战是：
Computer Use 技术（如 Anthropic 的 Claude Computer Use）
是否会在未来 2--3 年内成熟到可以替代传统 RPA 的程度？
如果 AI Agent 能够像人类一样``看屏幕、动鼠标''，
那么基于``预定义脚本模拟屏幕操作''的 RPA 就失去了存在的技术基础。
但至少在 2026 年的当下，
Computer Use 的准确性和可靠性还远未达到企业级要求——
它可以在 demo 中表现惊艳，
但在实际的企业环境中（复杂的 UI、多窗口切换、
权限弹窗、验证码、网络延迟），
错误率仍然过高，无法用于关键业务流程。
RPA 的确定性执行（deterministic execution）和
完整的审计日志（audit trail）
在可预见的未来仍然是企业客户的刚性需求。

\subsubsection{RPA 行业的未来：三大战略拐点}

Gartner 的 2025 年 RPA 魔力象限报告指出，
RPA 行业面临三大战略拐点：

\textbf{智能体自动化}（Agentic Automation）。
AI Agent 将成为 RPA 的协同引擎，
处理传统脚本无法解决的非确定性任务。
2024 年 Gartner 相关咨询量激增 750\%，
印证市场爆发的临界点已至。
未来，AI Agent 将动态调用 RPA 执行终端操作，
形成``决策--执行''闭环；
反之，RPA 也可触发 Agent 进行复杂推理。

\textbf{Computer Use}。
与传统 RPA 依赖预定义工作流不同，
Computer Use 技术（如 Anthropic 的 Claude Computer Use、
OpenAI 的 Operator）通过概率模型动态适应 UI 变化，
大幅降低维护成本。
尽管当前 Computer Use 尚无法完全替代 RPA
（准确性和可靠性仍有差距），
但其技术迭代速度可能在未来 2--3 年内重塑市场格局。
这是 UiPath 和 Automation Anywhere 面临的最大长期风险——
如果 AI Agent 能够直接``像人类一样操作电脑''，
传统 RPA 的``脚本化操作''范式就失去了存在价值。

\textbf{业务编排与自动化技术融合}（BOAT）。
RPA 正加速与智能文档处理（IDP）、低代码应用平台（LCAP）、
iPaaS（集成平台即服务）等技术栈融合，
形成统一的企业级自动化平台。
这一趋势模糊了 RPA、BPM（业务流程管理）、
iPaaS 和 AI Agent 平台之间的边界——
未来，这些类别可能会融合成一个统一的
``企业自动化 + AI''平台市场。

在进入垂直 SaaS 综合分析之前，图~\ref{fig:ent-heatmap}对本章讨论过的
十大行业的 AI 渗透程度进行了一次横向对比。

\begin{figure}[htbp]
\centering
\begin{tikzpicture}[
  rowlbl/.style={font=\footnotesize, anchor=east, text=black},
  cellH/.style={fill=figBlue!75, minimum width=2.8cm, minimum height=0.7cm,
    font=\scriptsize\color{white}, align=center},
  cellM/.style={fill=figBlue!40, minimum width=2.8cm, minimum height=0.7cm,
    font=\scriptsize, align=center},
  cellL/.style={fill=figBlue!15, minimum width=2.8cm, minimum height=0.7cm,
    font=\scriptsize, align=center},
  hdrlbl/.style={font=\footnotesize\bfseries, text=figGray, align=center},
]

% --- Column headers ---
\node[hdrlbl] at (3.0, 0.5) {Low};
\node[hdrlbl] at (5.8, 0.5) {Medium};
\node[hdrlbl] at (8.6, 0.5) {High};
\node[hdrlbl] at (0.2, 0.5) {Industry};

% --- Define rows: industry, penetration level, representative companies ---
% Row data: y-offset, name, level (H/M/L), company text
% High penetration
\foreach \y/\name/\comp in {
  0/Customer Service/Sierra{,} Intercom Fin,
  -0.8/Cybersecurity/CrowdStrike{,} Palo Alto,
  -1.6/Finance/AlphaSense{,} Ramp} {
  \node[rowlbl] at (1.1, \y) {\name};
  \node[cellL] at (3.0, \y) {};
  \node[cellM] at (5.8, \y) {};
  \node[cellH] at (8.6, \y) {\comp};
}

% Medium penetration
\foreach \y/\name/\comp in {
  -2.4/Retail/Shopify AI{,} Constructor,
  -3.2/Legal/Harvey{,} EvenUp,
  -4.0/Healthcare/Hippocratic{,} Abridge,
  -4.8/HR \& Recruiting/Eightfold{,} Paradox} {
  \node[rowlbl] at (1.1, \y) {\name};
  \node[cellL] at (3.0, \y) {};
  \node[cellM] at (5.8, \y) {\comp};
  \node[cellH] at (8.6, \y) {};
}

% Low penetration
\foreach \y/\name/\comp in {
  -5.6/Manufacturing/Sight Machine{,} Uptake,
  -6.4/Construction/Buildots{,} OpenSpace,
  -7.2/Agriculture/Aigen{,} Taranis} {
  \node[rowlbl] at (1.1, \y) {\name};
  \node[cellL] at (3.0, \y) {\comp};
  \node[cellM] at (5.8, \y) {};
  \node[cellH] at (8.6, \y) {};
}

% --- Grid lines ---
\foreach \y in {0.35,-0.45,-1.25,-2.05,-2.85,-3.65,-4.45,-5.25,-6.05,-6.85,-7.55} {
  \draw[figLightGray, line width=0.3pt] (1.6,\y) -- (10.0,\y);
}

% --- Legend ---
\node[cellL, minimum width=1.2cm] at (3.0, -8.4) {};
\node[font=\tiny, right] at (3.7, -8.4) {Low};
\node[cellM, minimum width=1.2cm] at (5.3, -8.4) {};
\node[font=\tiny, right] at (6.0, -8.4) {Medium};
\node[cellH, minimum width=1.2cm] at (7.6, -8.4) {};
\node[font=\tiny, right] at (8.3, -8.4) {High};
\node[font=\tiny\itshape, text=figGray] at (5.8, -9.0)
  {Penetration assessed by automation depth, market traction, and VC funding density};

\end{tikzpicture}
\caption{企业 AI 行业渗透率热力图：十大垂直领域对比（2025--2026）}
\label{fig:ent-heatmap}
\end{figure}

%% ============================================================
\section{垂直 SaaS + AI：全面颠覆}
\label{sec:ent-vertical-saas}
%% ============================================================

前面的章节按行业逐一分析了 AI 的渗透。
本节退后一步，从跨行业的视角审视一个更宏大的趋势：
\textbf{AI 正在系统性地重构每一个垂直 SaaS 市场}。

\subsection{``垂直 SaaS + AI''的三种范式}

AI 对垂直 SaaS 的改造呈现出三种不同的范式，
它们的颠覆程度递进，代表着企业 AI 的三个演化阶段：

\subsubsection{范式一：AI 增强（AI-Enhanced）}
\textbf{定义}：在现有的垂直 SaaS 产品中嵌入 AI 功能。
\textbf{示例}：Procore 在其建筑项目管理平台中加入 AI 风险预测；
ServiceNow 在其 ITSM 平台中加入 Now Assist；
Veeva 在其生命科学 CRM 中加入 AI 生成的临床文档摘要。
\textbf{特征}：AI 是``锦上添花''而非``雪中送炭''。
现有产品的核心逻辑不变，AI 作为附加功能提升用户体验。
去掉 AI 功能，产品依然完整可用。
\textbf{护城河}：依赖于原有产品的市场份额和客户粘性。
AI 功能容易被竞争对手复制（调用同一个 LLM API 即可）。

\subsubsection{范式二：AI 原生（AI-Native）}
\textbf{定义}：从第一天起就以 AI 为核心设计产品架构。
\textbf{示例}：Glean（AI 原生的企业搜索）；
ThoughtSpot（AI 原生的数据分析）；
Buildots（AI 原生的建筑监控）；
Vectra AI（AI 原生的网络检测）；
Landing AI（AI 原生的制造质检）。
\textbf{特征}：AI 不是``添加''的功能，而是产品存在的理由。
没有 AI，这些产品根本不会存在——
Buildots 离开计算机视觉就只是一堆没用的照片，
Vectra AI 离开行为分析就只是一个日志收集器。
\textbf{护城河}：行业特定的 AI 模型和数据飞轮。
每一次客户使用都在贡献训练数据，
使模型对这个行业的理解越来越深。

\subsubsection{范式三：AI Agent 重构（Agent-Restructured）}
\textbf{定义}：用 AI Agent 替代整个垂直 SaaS 的``人工服务''层。
\textbf{示例}：Salesforce Agentforce 替代人类客服 Agent；
Sierra AI 用 AI Agent 替代整个客户服务团队；
Harvey 用 AI Agent 替代初级律师的法律研究工作；
Abridge 用 AI Agent 替代医疗文档记录员。
\textbf{特征}：不是``让人类用更好的工具工作''，
而是``让 AI 直接做人类的工作''。
这是对垂直 SaaS 商业模式最深刻的颠覆——
它不仅改变了产品，还改变了定价逻辑
（从``按用户数''到``按工作量''或``按结果''）。
\textbf{护城河}：行业工作流的深度理解 + Agent 的执行可靠性 +
对错误的零容忍能力（在法律和医疗等领域，
Agent 的一次错误可能导致诉讼或生命风险）。

\subsection{典型的``范式迁移''案例}

为了更具体地说明这三种范式的差异，
让我们通过几个跨行业的案例来观察``范式迁移''是如何发生的。

\subsubsection{客户服务：从 Zendesk AI 到 Sierra AI}
Zendesk 是全球最大的客户服务 SaaS 平台之一，
2023 年在其产品中添加了 AI 功能——
自动分类工单、建议回复模板、生成客户摘要。
这是典型的``范式一：AI 增强''——
AI 功能是``锦上添花''，去掉它，Zendesk 依然是一个完整的客服平台。

Sierra AI 则代表了``范式三：AI Agent 重构''。
由 Salesforce 前联合 CEO Bret Taylor 和 Google AI 前负责人共同创立，
Sierra 的核心产品不是``帮助人类客服做得更好''，
而是``用 AI Agent 直接替代人类客服''。
Sierra 的 AI Agent 能够完全自主地处理客户的退货请求、
账户变更、技术支持等问题，
无需任何人类干预。
这种范式对 Zendesk 的``按座席收费''商业模式构成了根本性威胁——
如果企业可以用 Sierra 的 AI Agent 处理 80\% 的客户请求，
它就只需要 20\% 的人类客服座席，
对 Zendesk 的付费座席需求直接减少 80\%。

\subsubsection{法律科技：从 Westlaw AI 到 Harvey}
Thomson Reuters 的 Westlaw 是全球最大的法律数据库之一，
近年来添加了 AI 辅助的法律检索和摘要功能——
典型的``范式一''。
Harvey 则是法律 AI 领域的``范式二/三''代表——
一个 AI 原生的法律工作平台，
其 AI Agent 能够独立完成法律研究、合同审查、
文件起草等初级律师的核心工作。
Harvey 不是``让律师检索法规更快''，
而是``让 AI 直接做律师做的事''。

\subsubsection{医疗文档：从 Epic AI 到 Abridge}
Epic Systems 是美国最大的电子病历（EMR）系统供应商。
它在 EMR 中嵌入了 AI 功能——
辅助编码、药物交互提醒、智能排程——典型的``范式一''。
Abridge 则专注于用 AI Agent 自动生成医疗文档——
医生和患者的对话被 AI 实时转录、理解和结构化，
自动生成符合医疗规范的病历记录。
这不是``帮医生更快地写病历''，
而是``AI 直接写病历，医生只需审核''——
为每位医生每天节省 2--3 小时的文档工作时间。

这些案例说明了一个核心趋势：
\textbf{在每一个垂直领域，
都在同时发生``范式一→范式二→范式三''的迁移}。
传统的行业龙头在做``范式一''（AI 增强），
新兴的 AI 原生创业公司在做``范式二/三''（AI 原生/Agent 重构）。
两者之间的竞争将决定未来 3--5 年企业软件市场的格局重组。

\subsection{跨垂直领域的共性模式}

纵观本章讨论的所有垂直领域，
几个共性模式清晰浮现：

\textbf{模式一：数据飞轮效应}。
无论是 Glean 的企业知识图谱、CrowdStrike 的威胁情报、
Buildots 的建筑影像库、Taranis 的作物影像数据集，
还是 o9 Solutions 的供应链决策数据，
成功的垂直 AI 公司都在积累行业特定的数据资产。
这些数据越多，AI 模型越准确；
模型越准确，客户越愿意使用；
客户越多，数据越丰富——
这就是典型的数据飞轮（data flywheel）。
这种飞轮一旦启动，就会形成赢者通吃的竞争格局——
后来者即使有同等的技术能力，
也无法在短期内积累先行者已经拥有的数据资产。

\textbf{模式二：从``软件''到``决策''的价值迁移}。
传统 SaaS 的价值在于``记录''（System of Record）
和``流程''（System of Process）——
CRM 记录客户信息，ERP 管理业务流程，
项目管理工具追踪任务进度。
AI 将 SaaS 的价值迁移到``洞察''（System of Insight）
和``行动''（System of Action）——
不仅记录发生了什么，还告诉你应该做什么，
甚至直接帮你做了。
这种价值迁移使得 SaaS 公司能够从``按座席收费''
转向``按价值收费''——
如果 AI 帮你避免了 100 万美元的供应链中断损失，
收取 10 万美元的服务费就是理所当然的。
这从根本上改变了 SaaS 行业的单位经济学。

\textbf{模式三：平台 vs 垂直的永恒张力}。
本章反复出现的一个主题是横向平台与垂直专家之间的竞争与合作。
Microsoft、Google、Salesforce、ServiceNow 等平台巨头
拥有渠道、客户关系和数据优势，
但缺乏行业深度。
Glean、Buildots、Vectra AI、Landing AI 等垂直专家
拥有行业知识和精准的 AI 模型，
但缺乏平台规模和销售团队。
两者之间的关系不是简单的``竞争''或``互补''，
而是动态演化的——
今天的合作伙伴可能是明天的竞争对手
（Moveworks 曾是 ServiceNow 的集成伙伴，后被其收购）；
今天的竞争对手可能是明天的生态合作者
（UiPath 与 Microsoft Copilot Studio 的集成）。
垂直 AI 公司的生存策略是：
\textbf{要么做到行业第一，建立平台无法复制的数据壁垒；
要么找到平台的战略盲区，在夹缝中建立不可替代的价值。}

\textbf{模式四：地缘政治驱动的市场分裂}。
中美脱钩不仅影响芯片和模型层（参见第~\ref{ch:models}章），
也深刻影响着企业 AI 的应用层。
来也科技在中国信创市场的成功、
Altana AI 在贸易合规领域的崛起，
都是地缘政治重塑企业 AI 市场的体现。
未来，我们可能会看到两个相对独立的企业 AI 生态系统——
一个以美国技术栈为核心
（Microsoft + Salesforce + ServiceNow + UiPath + AWS/Azure），
一个以中国技术栈为核心
（来也科技 + 用友/金蝶 + 通义千问/文心 + 国产云平台）。
全球化的企业将不得不同时运营两套 AI 系统——
这是一个巨大的挑战，也是一个巨大的机会。

\subsection{长尾垂直：企业 AI 的千行百业全景补遗}

前面的章节聚焦于最大的几个赛道——
搜索、安全、供应链、制造、建筑、农业、RPA。
但企业 AI 的长尾远比这些``头部赛道''更加丰富。
事实上，每一个行业、每一个职能，
几乎都已诞生了专属的 AI 创业公司。
本节以``广度优先''的方式扫描这些被主流媒体忽略、
但在各自垂直领域中构建着深厚壁垒的``隐形冠军''。

\subsubsection{企业搜索与知识管理的长尾}

Glean 是企业搜索赛道的明星，
但这条赛道上还有大量差异化的竞争者：

\textbf{Hebbia}（纽约，2020 年成立，约 124 人）。
专注于金融行业的 AI 搜索和分析平台。
与 Glean 面向通用企业不同，Hebbia 深耕投资银行、私募信贷和对冲基金，
能够在数千份交易文档、PDF、备忘录和邮件中进行深度语义搜索。
其``Matrix''产品被投资专业人士视为``AI 版 Bloomberg Terminal''——
将原本需要数十小时的尽职调查压缩到分钟级。
2024 年获得 Andreessen Horowitz 领投的融资，
估值跻身独角兽行列。

\textbf{Vectara}（硅谷，2022 年成立，约 60 人）。
为开发者提供构建 AI 搜索和问答应用的平台。
核心技术是 RAG（检索增强生成），主打``零幻觉''——
通过精确的检索和引用追溯，最大化降低 LLM 的幻觉率。
定位为``开发者优先''的企业搜索基础设施，
与 Glean 的``开箱即用''路线形成差异化。

\textbf{Sinequa}（法国巴黎，2002 年成立，约 300 人）。
欧洲最大的企业 AI 搜索平台之一，
2025 年 Q4 被 QKS Group 评为 SPARK Matrix 企业 AI 搜索领导者。
专注于大型企业和政府机构的复杂搜索需求，
在航空航天（Airbus）、制药（Pfizer）和国防领域有深厚客户基础。
其``Enterprise Agentic AI''平台将搜索与 Agent 结合，
支持多步骤的知识发现和自动化工作流。

\textbf{Coveo}（加拿大魁北克，2005 年成立，约 700 人，TSX 上市）。
AI 搜索和推荐平台，主攻电商搜索、客户自助服务和企业知识发现。
与 Glean 的``内部搜索''定位不同，
Coveo 侧重于面向客户的搜索体验——
帮助企业客户更快地找到产品、解决方案和技术支持文档。
年收入约 1.3 亿加元。

\textbf{达观数据}（上海，2015 年成立，约 300 人）。
中国领先的企业级智能文本处理和知识管理公司。
以自主研发的``曹植''大语言模型为核心，
提供智能文档处理（IDP）、智能知识管理系统（KMS）、
RPA 和 AI Agent 平台等产品矩阵。
连续五年上榜毕马威金融科技双 50 榜单，
获吴文俊人工智能奖，拥有 100 余项发明专利。
典型客户覆盖银行、证券、制造等行业。

\textbf{Vespa.ai}（挪威，Yahoo 开源项目独立，2023 年成立）。
开源的 AI 搜索平台，支持向量搜索、张量计算和实时推理。
被 Spotify、Yahoo 等公司用于大规模搜索和推荐系统。
与 Glean 的 SaaS 模式不同，
Vespa 定位为搜索基础设施的开源替代方案——
面向有强大工程团队、需要深度定制的企业客户。

\subsubsection{AI 合规与法务}

法律科技是企业 AI 中最具``知识密集型''特征的赛道之一。
除了前文提到的 Harvey 和 Westlaw AI，
大量创业公司正在攻克合同审查、合规监控和法律研究的自动化。

\textbf{Luminance}（英国伦敦，2015 年成立，约 300 人）。
全球法律 AI 的标杆企业之一。
其``Legal-Grade AI''技术能够自主进行合同谈判——
2026 年 3 月发布的增强版是业界首个
能``展示每个决策背后原因''的全自主合同谈判工具。
2025 年推出 Compliance Agent，自动执行合规检查工作流。
核心创新是``机构记忆''（Institutional Memory）——
AI 系统能够学习一个法律团队的历史决策模式和偏好，
据称可为法律团队节省 30\% 的时间。
服务全球 500 多家企业和律所。

\textbf{ContractPodAi}（英国/美国，2012 年成立，约 200 人）。
AI 驱动的合同生命周期管理（CLM）平台。
2025 年 UNLEASHED 大会展示了其 Agentic AI 能力——
AI 不仅审查合同，还能自主起草、修订和管理整个合同生命周期。
被 Gartner 列为 CLM 市场的领导者。

\textbf{BlackBoiler}（美国，2017 年成立，约 30 人）。
聚焦于合同标记（Contract Markup）的自动化。
AI 能够按照企业预设的法律标准和语言风格，
自动对合同进行逐条审查和修改建议。
获多项法律 AI 奖项，
定位比 Luminance 更``窄''但更``深''——
只做合同审查这一件事，做到极致。

\textbf{Sirion}（原 Eigen Technologies，美国/英国，2012 年成立，约 785 人）。
2025 年 Eigen Technologies 被 Sirion 收购并整合，
形成了 AI 合同管理领域的大型平台。
Sirion 年收入约 1.6 亿美元，结合了 Eigen 的智能文档处理（IDP）技术
和自身的合同生命周期管理能力，
推出了``对话式 AI''合同管理界面。

\textbf{幂律智能}（北京，2018 年成立）。
中国法律 AI 的代表性创业公司。
2023 年完成近 8000 万元 Pre-B 轮融资，
蓝驰创投领投，电子签厂商 e 签宝跟投。
专注于法律大模型和合同智能审查，
是国内少数获得主流 VC 持续加注的法律 AI 公司。

\subsubsection{AI 财务与会计}

财务是企业中最``规则密集''的职能之一——
月结、对账、发票处理、预算编制、审计准备——
大量重复性但又需要专业判断的工作，
是 AI 自动化的天然土壤。

\textbf{Truewind}（硅谷，2022 年成立，约 50 人）。
AI 驱动的会计自动化平台，
能够自动完成月末结账中最繁琐的环节——
交易分类、银行对账、预付费摊销和波动分析。
2025 年入选 AICPA 和 CPA.com 创业加速器，
服务 500 多家会计师事务所和创业公司。
被 Tooliverse 评为 2026 年 AI 会计自动化 9.38/10 分。

\textbf{Vic.ai}（美国/挪威，2017 年成立，约 200 人）。
专注于应付账款（AP）自动化的 AI 平台。
其 AI 基于超过 10 亿张发票训练，
能够自动处理发票录入、审批路由和支付执行。
据 Vic.ai 自己的调查显示，
2025 年已有 72\% 的 AP 或财务团队在使用某种形式的 AI。
被视为财务 AI 领域最专业的``纯 AP''解决方案。

\textbf{Stampli}（硅谷，2015 年成立，约 250 人）。
AP 自动化平台，独特之处在于``以发票为中心''的协作界面——
将沟通、审批和 AI 建议集中在每张发票的上下文中。
AI 能力包括智能编码、异常检测和自动三方匹配。
与 Vic.ai 的``全自动''路线不同，
Stampli 更强调``人机协作''——
AI 提供建议，人类做最终决策。

\textbf{Tipalti}（硅谷，2010 年成立，约 1300 人）。
全球支付和应付账款自动化平台，
2025 年推出 Tipalti AI Assistant 和 AI Agents，
将 AI 嵌入从发票处理到跨境支付的全流程。
2024 年估值约 83 亿美元，年收入超 3 亿美元。
定位比 Vic.ai 和 Stampli 更``广''——
不仅做 AP 自动化，还覆盖供应商管理、合规和全球支付。

\subsubsection{AI 物流与仓储机器人}

在供应链的``最后一英里''和仓储环节，
AI + 机器人的融合正在创造全新的自动化范式。

\textbf{Locus Robotics}（马萨诸塞州，2014 年成立，约 600 人）。
仓储协作机器人（AMR）领域的领导者。
2025 年推出 LocusONE 平台——
将``Physical AI''（物理世界的 AI）和仓储领域知识融合，
实现 Agentic AI 驱动的仓储自动化。
旗舰产品 Locus Array 是业界首个``全机器人''履约系统，
实现零人工触碰的端到端订单处理。
与 DSV 等全球物流巨头建立了深度合作。

\textbf{Covariant}（加州伯克利，2017 年成立，约 150 人）。
AI 仓储机器人公司，核心技术是``Covariant Brain''——
一个让机器人能够理解和操作各种形状、材质物品的 AI 系统。
2024 年被 Amazon 收购了部分团队和技术，
创始人 Pieter Abbeel 是加州大学伯克利分校的机器人学教授。
代表了``从学术到工业''的 AI 机器人路径。

\textbf{Nuro}（硅谷，2016 年成立，约 1200 人）。
自动驾驶配送机器人公司，
由 Google 自动驾驶项目前首席工程师 Dave Ferguson 和 Jiajun Zhu 创立。
总融资超过 21 亿美元，估值约 86 亿美元。
其无人配送车已在多个美国城市商业运营，
与 7-Eleven、Domino's、FedEx 等品牌合作。
定位于``最后一英里''的完全自主配送。

\textbf{灵动科技}（北京，2016 年成立，约 400 人）。
中国仓储 AMR（自主移动机器人）的代表企业，
产品线覆盖从轻载分拣到重载搬运的全场景。
核心技术包括视觉 SLAM 导航和多机协同调度算法。
客户包括顺丰、京东物流、TCL 等，
也是少数成功出海的中国仓储机器人公司之一。

\subsubsection{AI 医疗（企业端）}

前文已在``垂直 SaaS + AI''中提及 Abridge（医疗文档）和 Epic AI，
但医疗 AI 的企业端远比这丰富——
从精准医疗到病理诊断，
从临床影像到医疗设备，
每个细分领域都有深耕的 AI 公司。

\textbf{Tempus AI}（芝加哥，2015 年成立，约 2200 人，NASDAQ: TEM）。
全球最大的 AI 精准医疗平台，年收入约 11 亿美元。
构建了世界最大的临床和分子数据库，
覆盖基因组学、蛋白质组学、影像学和电子健康记录。
其 AI 平台帮助肿瘤科医生做出更精准的治疗决策——
从``对症治疗''升级到``对因治疗''。
2024 年 IPO，是 AI 医疗领域最大的独立上市公司之一。

\textbf{PathAI}（波士顿，2016 年成立，约 200 人，收入约 2900 万美元）。
AI 驱动的数字病理学平台。
通过计算机视觉分析病理切片，
辅助病理学家做出更快、更准确的诊断——
尤其在肝脏活检和肿瘤分型领域。
为制药公司的临床试验提供 AI 病理分析服务，
同时运营 PathAI Diagnostics（孟菲斯，约 53 人）子公司。

\textbf{Viz.ai}（旧金山，2016 年成立，约 253 人，收入约 4200 万美元）。
AI 驱动的临床护理协调平台。
最初聚焦于中风检测——
AI 自动分析 CT 影像，在秒级内识别大血管闭塞，
并即时通知神经介入团队——将救治时间缩短数十分钟。
2025 年扩展到心脏病、肺栓塞等多个领域，
宣布年度达到创纪录的规模和患者影响力。

\textbf{Butterfly Network}（纽约，2011 年成立，NASDAQ: BFLY）。
将超声设备从数万美元的大型机器缩小到一个手持探头（Butterfly iQ+），
通过 AI 辅助图像解读，
让非专业人员也能进行基本的超声检查。
AI 自动识别器官、测量指标并提供诊断建议。
被称为``超声界的 iPhone''——
用硬件+AI 的组合颠覆传统医疗设备市场。

\textbf{推想科技}（北京，2016 年成立，约 500 人）。
中国 AI 医疗影像的头部企业之一，
以``AI 改善生命健康''为愿景，
采取``一横一纵''和国际化战略布局。
产品覆盖肺结节筛查、骨折检测、脑卒中预警等场景，
2025 年有望实现盈利——
这对整个中国 AI 医疗影像行业来说是一个里程碑信号。

\textbf{数坤科技}（北京，2017 年成立，约 500 人）。
AI 医疗影像独角兽，2021 年估值一度高达 94 亿元。
自主研发``数字心''``数字脑''``数字胸''``数字腹''等
``数字医生''产品系列，
覆盖心脏病、脑卒中、癌症等重大疾病的智能诊疗。
截至 2025 年已获批 19 张三类医疗器械注册证，
为行业持证数量最多的企业。
2025 年已启动 IPO 准备工作。

\textbf{鹰瞳科技}（北京，2015 年成立，港交所: 2251.HK）。
全球首批通过 AI 视网膜影像分析
进行慢性病筛查的医疗 AI 公司之一。
通过拍摄眼底照片，
AI 可以检测糖尿病视网膜病变、青光眼、高血压等多种疾病。
2021 年在港交所上市，
是中国``AI + 医疗器械''赛道少数成功 IPO 的企业。

\subsubsection{AI 教育（企业端）}

教育科技的企业端——面向学校、大学和培训机构的 AI 解决方案——
正在经历一场由 LLM 驱动的范式变革。

\textbf{Cengage Group}（美国，教育出版巨头）。
2025 年全面扩展 AI 产品线，
为数百万高等教育、K-12 和职业培训学习者提供个性化体验。
与 AWS 合作定义教育 AI 的未来，
推出了面向企业级部署的``Sharpen Advantage''平台。
代表了传统教育巨头``范式一：AI 增强''的路径。

\textbf{McGraw Hill}（美国，教育出版巨头，2025 年 IPO）。
推出 AI 驱动的``Sharpen''学习应用，
为学生提供个性化的学习路径和即时反馈。
CEO 在 2025 年 IPO 路演中将 AI 定位为公司未来增长的核心引擎。
代表了教育出版行业``All-in AI''的战略转型。

\textbf{松鼠 AI}（上海，2014 年成立）。
中国自适应学习 AI 的先行者，
核心技术是基于知识图谱的``AI 老师''——
能够精准诊断每个学生的知识薄弱点，
动态生成个性化学习路径。
曾获数亿元融资，服务数千家教育机构。
在``双减''政策冲击后，
转型面向成人教育和企业培训市场。

\subsubsection{AI 零售}

零售业的线下门店正在成为计算机视觉和 AI 的主战场——
从防损、货架监控到无人结账，
AI 正在重新定义``智慧零售''。

\textbf{Everseen}（爱尔兰科克，2007 年成立，约 394 人，收入约 9300 万美元）。
零售视觉 AI 的全球领导者。
2026 年 1 月推出 Everact——
业界首个``对话式零售智能''Agentic AI 平台，
将门店视觉 AI 与自然语言交互结合。
旗下产品覆盖结账防损（Evercheck）、
货架监控（Evershelf）和批量结账优化等场景。
与 Google Cloud 深度合作，在 12 个国家运营。

\textbf{Standard AI}（旧金山，2017 年成立，约 100 人）。
计算机视觉零售平台，
利用 AI 分析门店内的购物行为和商品状态——
从货架缺货预警到顾客动线分析。
2025 年 NRF 大会上展示了``计算机视觉洞察''能力，
为零售商和消费品品牌提供可操作的门店智能。

\textbf{Trax}（新加坡，2010 年成立，约 500 人）。
AI 驱动的零售图像识别平台，
号称 96\% 的商品识别准确率。
通过手机拍照、固定摄像头或机器人巡检，
AI 实时分析货架陈列、价格标签和促销执行情况。
服务于宝洁、联合利华等全球消费品巨头。

\textbf{影刀 RPA}（杭州，2019 年成立，约 300 人）。
虽然名为``RPA''，但影刀已演化为面向电商和零售企业的
AI 自动化全平台。
覆盖从商品上架、订单处理、客服自动回复到数据报表的全流程，
在淘宝、京东、拼多多等电商生态中被广泛使用。
用户数超过百万，是中国 RPA 领域用户量最大的产品之一。

\subsubsection{AI 建筑与房地产（补充）}

前文已详细讨论了 OpenSpace 和 Buildots，
这里补充几家更``小众''但在各自领域有独特价值的公司。

\textbf{TestFit}（达拉斯，2015 年成立，约 60 人）。
房地产可行性分析的 AI 平台。
其``生成式设计''引擎能够在数小时内
生成数百种建筑方案并进行财务可行性分析——
传统方法需要数天甚至数周。
2025 年入选 BuiltWorlds 建筑科技 Top 50，
支持多户住宅、数据中心、工业园区等多种建筑类型。

\textbf{Swapp}（以色列特拉维夫，2018 年成立，约 60 人）。
AI 建筑文档自动化平台。
将建筑师的概念设计自动转化为完整的施工文档和 BIM 模型——
包括结构计算、MEP（机电管道）布局和合规检查。
号称``按照建筑事务所自己的标准和偏好''进行自动化，
将数月的文档工作压缩到数天。

\textbf{小库科技}（深圳，2016 年成立，约 100 人）。
中国领先的 AI 建筑设计平台，
核心产品是``AI 强排''——
通过生成式 AI 自动生成住宅小区的建筑排布方案，
优化容积率、日照、通风和景观等多维指标。
服务中国头部房企和建筑设计院，
是中国``AI + 建筑''赛道的先行者。

\subsubsection{AI 人力资源}

HR 是企业中``人''最密集的职能——
招聘、绩效、薪酬、培训、员工发展——
每个环节都涉及大量的判断、沟通和文档工作。
AI 正在系统性地渗透 HR 的每一个流程节点。

\textbf{Phenom}（美国费城，2011 年成立，约 1500 人）。
``Applied AI''人才智能平台，
覆盖从招聘、入职到留存的全生命周期。
核心能力是将非结构化的``人才数据''（简历、面试记录、
绩效评估、技能标签）转化为可操作的智能洞察。
Phenom 不只是一个招聘工具，
而是一个``人才操作系统''——
帮助企业理解``我们需要什么样的人''和``我们有什么样的人''之间的差距。

\textbf{Beamery}（英国伦敦，2014 年成立，约 300 人）。
AI 驱动的``劳动力转型平台''，
将人才战略与业务成果连接。
2025 年发布了涵盖 AI 在 HR 中应用的深度白皮书，
核心观点是 AI 正在从``工作定义''到``任务执行''全面重塑组织——
不仅改变``如何招人''，还改变``如何定义工作本身''。

\textbf{Findem}（硅谷红木城，2019 年成立，约 168 人，收入约 3900 万美元）。
AI 人才获取和劳动力决策平台。
2025 年完成 3600 万美元 C 轮融资。
核心差异化在于``深度上下文''——
不仅匹配技能和经验，
还理解候选人的职业轨迹、公司文化适配度和隐性潜力。
从招聘扩展到内部流动、学习发展和劳动力规划。

\textbf{SeekOut}（西雅图，2017 年成立，约 200 人）。
AI 人才搜索和多元化招聘平台。
核心技术是跨平台的``人才图谱''——
整合 LinkedIn、GitHub、专利数据库等多源数据，
构建每个候选人的 360 度画像。
特别强调多元化和包容性（DEI）——
帮助企业在招聘中系统性地减少偏见。

\textbf{北森}（北京，2002 年成立，约 3000 人）。
中国最大的一体化 HR SaaS 平台，
2022 年在港交所上市（6651.HK）。
2025 年推出``AI Family''全产品矩阵——
覆盖``选、育、用、留''全生命周期，
连续九年蝉联 IDC 中国 HR SaaS 市场份额第一。
2025 年 HR SaaS 领域 AI Agent 产品 TOP10 中，
北森 AI Family 的商业化进度遥遥领先。

\textbf{Moka}（北京，2015 年成立，约 400 人）。
中国新一代 AI 招聘和人力资源管理系统。
2025 年推出``Moka Eva''AI 招聘智能体——
从``招聘助手''升级为能够自主执行简历筛选、
面试安排和候选人沟通的 AI Agent。
被视为``AI 原生''HR SaaS 的代表——
与北森``传统 HCM + AI 增强''的路径形成差异化。

\subsubsection{AI 客户服务}

前文已在``范式迁移''中讨论了 Sierra AI 和 Zendesk AI。
客户服务是企业 AI 渗透最快的领域之一——
Gartner 预测到 2029 年，
Agentic AI 将自主解决约 80\% 的客户服务问题。

\textbf{Ada}（加拿大多伦多，2016 年成立，约 260 人，收入约 3100 万美元）。
AI 客户服务自动化平台，
宣称其 AI Agent 能够解决超过 70\% 的客户服务请求。
2025 年员工增长 63.5\%，是增长最快的客服 AI 公司之一。
与 Sierra 的``面向大品牌''定位不同，
Ada 更面向中型企业和成长型公司。

\textbf{Forethought}（旧金山，2018 年成立，约 100 人）。
多 Agent 客户支持 AI 平台，
覆盖电商、零售、金融科技和医疗等多个垂直领域。
核心架构是``多 Agent 系统''——
不同的 AI Agent 负责不同类型的客户问题，
协同完成复杂的客户请求。
定位比 Ada 更``平台化''，
更强调跨领域的通用客服 AI 能力。

\textbf{Intercom}（旧金山，2011 年成立，约 1023 人，收入约 2 亿美元）。
``AI 优先''的客户服务平台，
将 AI Agent、AI Copilot、工单系统和电话支持整合在一个平台中。
2025 年员工增长 23.4\%。
与专注于``AI 替代人类客服''的 Sierra 和 Ada 不同，
Intercom 走的是``AI + 人类协作''的路线——
AI 处理简单问题，复杂问题无缝转交给人类客服。

\subsubsection{AI 销售与营收智能}

企业销售的核心痛点是``信息不对称''——
销售代表不知道该联系谁、说什么、何时跟进；
销售管理者不知道哪些交易会成交、哪些会流失。
AI 正在从根本上解决这些信息不对称。

\textbf{Gong}（硅谷帕洛阿尔托，2015 年成立，约 1500 人）。
销售对话分析的开创者，
通过 AI 分析每一通销售电话、每一封邮件、每一次视频会议。
2025 年发布了``年度最佳销售洞察''，
基于数百万次真实销售对话的数据分析。
核心价值是让销售团队``用数据而非直觉做决策''。
估值约 73 亿美元，是 AI 销售赛道最大的独立公司。

\textbf{Clari}（硅谷桑尼维尔，2012 年成立，约 600 人）。
营收智能平台，2025 年被 Gartner 评为
首届``Revenue Action Orchestration''魔力象限的领导者。
2025 年宣布与 Salesloft 合并——
两家公司结合形成``预测性营收系统''（Predictive Revenue System），
将 AI 驱动的营收预测与销售执行合为一体。
这是 AI 销售赛道最大的并购事件之一。

\textbf{Outreach}（西雅图，2014 年成立，约 900 人）。
AI 营收工作流平台，
号称是``唯一统一了潜客开发、客户留存和营收智能''的 AI 平台。
2025 年推出 AI Revenue Agents——
AI 能够自动完成潜客研究、邮件个性化和跟进排程。
定位为销售团队的``AI 操作系统''。

\textbf{6sense}（旧金山，2013 年成立，约 700 人）。
AI 驱动的``意图数据''平台，
核心能力是识别``正在市场中寻找解决方案但尚未主动联系你''的潜在买家。
通过分析海量的网页行为、搜索数据和内容消费模式，
6sense 能够在买家主动出击之前就识别其购买意图——
被称为``暗漏斗''（dark funnel）的照明者。

\subsubsection{AI 采购与支出管理}

采购是企业中最容易``漏钱''的环节——
非合规采购、价格不透明、供应商风险未知。
AI 正在让采购从``成本中心''变成``价值中心''。

\textbf{Zip}（旧金山，2020 年成立，约 400 人）。
全球领先的 AI 采购编排平台，
2023 年以来已交付 570 万条 AI 洞察。
核心理念是``让采购流程对组织中的每一个人都易于使用''——
通过 AI Agent 自动化从需求发起到采购执行的全流程。
与 Fairmarkit 建立战略合作，
推出``自主 Source-to-Pay''能力。

\textbf{Coupa}（硅谷，2006 年成立，约 3500 人，2023 年被 Thoma Bravo 以 80 亿美元收购）。
``AI 原生''的全面支出管理平台，
覆盖从采购、发票到供应链和财资管理。
凭借全球最大的``Business Spend Network''（商业支出网络），
Coupa 的 AI 基于数万亿美元的交易数据进行训练——
数据飞轮效应极为显著。

\textbf{JAGGAER}（美国/德国，1995 年成立，约 1200 人）。
2025 年推出``JAI''——
一个``采购原生''的 AI Agent 编排器。
不同于通用 AI 助手，
JAI 专门针对采购和供应链流程进行训练，
能够自主完成供应商寻源、价格谈判辅助和合规检查。

\textbf{Fairmarkit}（波士顿，2017 年成立，约 100 人）。
``自主寻源''（Autonomous Sourcing）AI 平台。
AI 能够自动发出询价、分析报价、推荐最佳供应商——
将原本需要采购人员数天的寻源流程压缩到数分钟。
2025 年获 Coupa 认证，2026 年与 Zip 建立合作。

\subsubsection{AI 保险}

保险是一个``精算驱动''的行业——
核心业务（定价、理赔、风控）本质上是数据分析问题。
AI 正在将保险从``规则驱动''升级为``学习驱动''。

\textbf{Lemonade}（纽约，2015 年成立，约 1500 人，NYSE: LMND）。
从第一天起就以 AI 为核心的保险公司——
不是``用 AI 辅助保险''，而是``AI 就是保险''。
AI Maya 负责销售（秒级完成报价），
AI Jim 负责理赔（秒级完成理赔审核）。
2025 年成为宠物保险领域增长最快的公司，
全年业绩强劲增长，
证明了``AI 原生保险''模式的商业可行性。

\textbf{Tractable}（英国伦敦，2014 年成立，约 300 人，独角兽）。
AI 驱动的保险理赔自动化。
通过计算机视觉分析事故照片，
AI 能够在秒级内评估车辆损伤程度和维修成本。
2022 年完成 D 轮 6000 万美元融资后达到独角兽估值。
服务全球多家顶级保险公司和汽车制造商。

\textbf{Shift Technology}（法国巴黎，2014 年成立，约 400 人）。
保险行业领先的 AI 平台，
2025 年推出``Shift Claims''——
用 Agentic AI 驱动理赔全流程自动化。
核心能力是欺诈检测——
AI 从数百万历史理赔中学习欺诈模式，
帮助保险公司减少理赔漏损（claims leakage）。
在欧洲保险市场有极高渗透率。

\subsubsection{AI 药物发现}

AI 对制药行业的最大承诺是``10 倍 R\&D 效率''——
将传统需要 10 年、耗资 25 亿美元的新药研发周期
大幅压缩。

\textbf{Recursion Pharmaceuticals}（盐湖城，2013 年成立，约 1000 人，NASDAQ: RXRX）。
2024 年与 Exscientia 合并，
形成了 AI 药物发现领域最大的公司。
核心技术是``生物学地图''——
通过高通量生物实验+AI 分析，
系统性地探索疾病和药物的相互作用。
合并后的管线覆盖肿瘤、神经和罕见病等领域。

\textbf{Insilico Medicine}（香港/美国，2014 年成立，约 400 人）。
端到端的 AI 药物发现平台，
拥有``Chemistry42''（分子生成）、``Biology42''（靶点发现）
和``Clinical42''（临床试验预测）三大 AI 引擎。
2025 年其 AI 发现的抗纤维化药物 INS018\_055 进入 II 期临床——
是全球首批由 AI 从``靶点发现''到``候选分子''
全程设计的药物之一。

\textbf{Relay Therapeutics}（波士顿，2016 年成立，约 400 人，NASDAQ: RLAY）。
利用``动力学药物发现''（dynamo platform）——
结合 AI 和物理模拟来理解蛋白质运动，
从而设计更精准的小分子药物。
2025 年多个管线进入临床阶段。
代表了``AI + 物理模拟''的药物发现路径。

\subsubsection{AI 文档处理（IDP）}

智能文档处理（Intelligent Document Processing, IDP）
是企业 AI 中最``接地气''的应用之一——
将纸质文件、PDF、扫描件中的信息提取并结构化。

\textbf{ABBYY}（美国/德国，1989 年成立，约 1000 人）。
2025 年被 Gartner 评为 IDP 魔力象限领导者。
从传统 OCR 公司转型为 AI 驱动的 IDP 平台，
结合文档 AI 和 GenAI 的混合架构。
在银行、保险和物流等行业有极广的客户基础。

\textbf{Hyperscience}（纽约，2014 年成立，约 182 人）。
2025 年获``AI Breakthrough''年度 IDP 平台大奖。
核心差异化在于``人在回路''（Human-in-the-Loop）架构——
AI 自动处理高置信度的文档，
低置信度的自动路由给人类审核，
审核结果又反馈给 AI 学习。
与 ABBYY 和 UiPath 共同位列 Gartner MQ 第一梯队。

\textbf{Rossum}（英国/捷克，2017 年成立，约 156 人）。
专注于发票和财务文档的 AI 处理。
将 IDP 深度整合到 AP 工作流中——
不仅``读懂''发票，还能直接触发审批和支付。
年收入约 400 万美元（处于快速增长阶段，员工同比增长 15.8\%）。

\subsubsection{AI 合规与治理（GRC）}

随着 AI 的广泛部署，
``如何治理 AI 本身''成为了一个新的企业需求。
AI 治理和合规（Governance, Risk, Compliance）
正在形成一个独立的赛道。

\textbf{OneTrust}（亚特兰大，2016 年成立，约 2500 人，估值约 45 亿美元）。
全球最大的隐私和数据治理平台，
2025 年全面转向``AI 治理''——
帮助企业管理 AI 模型的风险、偏见和合规性。
核心产品包括 AI 清单管理、模型风险评估和负责任 AI 框架。
被 Gartner 列为隐私管理领导者。

\textbf{Drata}（圣地亚哥，2020 年成立，约 500 人）。
合规自动化平台，支持 SOC 2、ISO 27001、HIPAA 等
数十种合规框架的持续监控和自动报告。
2025 年新增 NIST AI 风险管理框架支持——
帮助企业满足日益严格的 AI 合规要求。

\textbf{Vanta}（旧金山，2018 年成立，约 500 人，估值约 25 亿美元）。
与 Drata 类似的合规自动化平台，
被 G2 评为 2025 年 GRC 赛道的用户满意度第一。
2026 年发布``5 大最佳 GRC 解决方案''报告，
将自身定位为``AI 时代的合规基础设施''。

\textbf{LogicGate}（芝加哥，2015 年成立，约 200 人）。
2026 年被 Gartner 评为 GRC 魔力象限领导者。
其``Risk Cloud''平台以``有意图的 AI''
（Intentional AI）为设计原则——
确保 AI 在合规决策中的可解释性和可审计性。

\subsubsection{AI 网络安全（补充）}

前文已详细讨论了 CrowdStrike、SentinelOne 和 Vectra AI，
但网络安全 AI 赛道还有大量值得关注的创业公司。

\textbf{Abnormal Security}（旧金山，2018 年成立，约 550 人）。
AI 驱动的邮件安全和社会工程攻击防护平台。
核心技术是``行为基线''——
AI 学习每个用户和组织的正常通信模式，
检测偏离基线的异常行为（如 CEO 欺诈、供应商邮件劫持）。
被 Gartner 列为邮件安全的领导者。

\textbf{Wiz}（纽约/以色列，2020 年成立，约 1500 人，估值超 120 亿美元）。
2025 年被 Google Cloud 收购——
这是云安全赛道有史以来最大的并购之一。
核心产品是无 Agent 的云安全平台，
通过 AI 自动发现、优先排序和修复云环境中的安全风险。
从代码到运行时形成统一的安全图谱。

\textbf{Orca Security}（以色列特拉维夫，2019 年成立，约 500 人）。
与 Wiz 的直接竞争对手，
同样采用无 Agent 架构的云安全平台。
强调``深度优先''的安全分析——
不仅发现风险，还提供详细的攻击路径分析和修复指导。

\textbf{Snyk}（波士顿/伦敦，2015 年成立，约 1200 人，估值约 74 亿美元）。
AI 开发者安全平台，2025 年推出``AI Security Fabric''——
不仅保护代码安全，还保护 AI 模型和 Agent 的安全。
2025 年推出 Evo by Snyk——
专门为 Agentic AI 应用提供安全编排，
在 AI SDLC 全流程中防御新型威胁。
代表了``安全左移''和``AI 安全''的融合趋势。

\textbf{Prophet Security}（门洛帕克，2024 年成立，约 55 人，员工同比增长 230\%）。
AI SOC 分析师——
自主执行 Tier 1/Tier 2 安全告警的分类和调查。
2024 年刚成立就实现了爆发式增长，
代表了``AI Agent 替代安全分析师''的新一代安全创业公司。

\subsubsection{AI 可观测性与 IT 运维}

当企业 IT 环境变得越来越复杂
（微服务、容器、多云、AI 工作负载），
``可观测性''（Observability）成为关键需求——
AI 正在将可观测性从``仪表盘监控''升级为``自主诊断''。

\textbf{Dynatrace}（美国，2005 年成立，约 5000 人，NYSE: DT）。
2025 年将 AI 可观测性定位为核心战略——
不仅监控传统 IT 系统，
还为 LLM 和 Agentic AI 工作流提供可观测性，
包括性能优化、可解释性和合规监控。

\textbf{Chronosphere}（纽约，2019 年成立，约 200 人）。
云原生可观测性平台，
核心差异化在于``成本控制''——
帮助企业在海量遥测数据中
只存储和分析真正重要的数据，
将可观测性成本降低 60\% 以上。

\textbf{Observe, Inc.}（旧金山，2017 年成立，约 102 人，收入约 5400 万美元）。
AI 驱动的可观测性平台，
号称``故障排查速度提升 10 倍，成本降低 60\%''。
虽然员工同比下降 46.1\%（经历了行业调整期），
但年收入仍保持稳健。

\subsubsection{中国企业 AI 生态补遗}

除了前文已讨论的来也科技、达观数据、北森和 Moka，
中国企业 AI 生态还有更多值得关注的玩家：

\textbf{飞书 AI}（字节跳动，北京）。
字节跳动的企业协作平台``飞书''深度整合了 AI 功能——
AI 会议纪要、智能文档助手、自动化工作流。
2025 年推出``飞书智能伙伴''，
基于字节的豆包大模型，
提供从会议总结到数据分析的全场景 AI 辅助。
凭借字节的技术积累和中国市场的规模优势，
飞书 AI 正成为中国企业协作赛道的``AI 标杆''。

\textbf{钉钉 AI}（阿里巴巴，杭州）。
中国最大的企业协作平台之一，
2025 年全面接入通义千问大模型，
推出智能助理、智能会议、AI 审批和智能报表等功能。
依托阿里云的基础设施和阿里生态的客户资源，
钉钉 AI 在制造业、零售和中小企业市场有极高渗透率。

\textbf{企业微信 AI}（腾讯，深圳）。
基于微信生态的企业协作平台，
2025 年接入腾讯混元大模型。
核心优势是``连接微信''——
企业员工可以直接在企业微信中与 12 亿微信用户沟通，
AI 辅助客户管理、营销自动化和内部协作。
在零售、教育和医疗行业有独特的生态优势。

\textbf{金山办公 AI}（WPS AI，珠海，深交所: 688111）。
WPS Office 的 AI 升级版，
通过 AI 实现文档生成、PPT 设计、表格分析和 PDF 处理。
月活用户超过 5 亿，
是全球用户量最大的办公 AI 产品之一（在中国市场）。
依托庞大的用户基础，
WPS AI 正在从``工具''进化为``AI 工作平台''。

\subsection{2026--2028 年展望}

展望未来两到三年，企业 AI 赛道将呈现以下趋势：

\begin{enumerate}[noitemsep]
  \item \textbf{Agent 化全面展开}。
        到 2028 年，大多数企业 SaaS 产品将内置 AI Agent 能力。
        不具备 Agent 功能的 SaaS 产品将面临客户流失和市场份额萎缩。
        ``你们的产品有 Agent 吗？''将成为企业采购的标准问题。
  \item \textbf{定价模式重构}。
        按``对话数''``Agent 执行数''``业务结果''收费的模式
        将逐步取代传统的按座席收费。
        这将重塑整个 SaaS 行业的单位经济学——
        一些公司的 ARPU 会因此上升（如 Salesforce），
        另一些公司的 ARPU 会因此下降（如传统按人头收费的 SaaS）。
  \item \textbf{整合浪潮加速}。
        像 ServiceNow 收购 Moveworks（28.5 亿美元）、
        Automation Anywhere 收购 Aisera 这样的整合将加速。
        独立的``点状''AI 公司如果不能成为平台，
        就会被平台收购。
        2026--2028 年可能是企业 AI 赛道的``并购黄金期''。
  \item \textbf{行业 AI 标准化}。
        各垂直行业将逐步形成 AI 应用的标准框架和最佳实践。
        建筑行业的``AI 施工监控标准''、
        网络安全的``Agentic SOC 参考架构''、
        金融行业的``AI 合规审计规范''等将陆续出台。
        这些标准的制定者将获得显著的先发优势。
  \item \textbf{从自动化到自主化}。
        企业 AI 的最终愿景不是``更高效的人机协作''，
        而是``AI 自主执行关键业务流程''。
        从 RPA 的``模拟人类操作''到 Agentic AI 的``自主完成任务''，
        企业正在走向``自主企业''（Autonomous Enterprise）的终极形态——
        财务对账自动完成、客户工单自动处理、
        供应链异常自动预警并自动调整、安全威胁自动检测并自动响应。
        当然，这一愿景的实现需要解决信任、治理、合规和安全等
        一系列非技术问题——
        这些问题可能比技术问题本身更加棘手。
\end{enumerate}

\begin{warnbox}[企业 AI 的``最后一公里''挑战]
本章聚焦于成功案例和增长数据，
但必须指出，企业 AI 的实际部署远比这些光鲜的数字所暗示的更加艰难。
据 Andreessen Horowitz 对企业 CIO 的调查显示，
即使在最先进的企业中，LLM 的 ROI 也``并不如预期那样戏剧性''。
企业 AI 的``最后一公里''挑战包括：

\begin{itemize}[noitemsep]
  \item \textbf{数据质量}：AI 模型的输出质量取决于输入数据的质量。
        许多企业的内部数据是分散的、不一致的、过时的。
        ``垃圾进、垃圾出''（Garbage In, Garbage Out）
        在 AI 时代被放大了——
        传统软件给你一个错误的报表，你知道它是错的；
        AI 给你一个``看起来很对''但实际上基于错误数据的分析，
        你可能会信以为真并据此做出决策。
  \item \textbf{变更管理}：员工对 AI 的接受度是部署成功的关键。
        ``技术可行''不等于``组织接受''。
        ThoughtSpot 的研究显示，
        即使在``领先型''组织中，
        也只有 53\% 的员工能够即时获取 AI 洞察——
        意味着近一半的员工仍未被 AI 触达。
  \item \textbf{合规与审计}：在金融、医疗、政府等受监管行业，
        AI 的决策需要可解释、可审计、可追溯。
        一个 AI Agent 自主批准了一笔贷款——
        监管机构会问：``为什么？依据是什么？
        如果决策错误，谁负责？''
  \item \textbf{集成复杂度}：企业 IT 环境是由数十甚至数百个系统组成的
        ``意大利面条架构''（spaghetti architecture）。
        将 AI 集成到这样的环境中，
        工程挑战往往远超 AI 模型本身——
        一个 AI 助手需要连接 15 个不同的系统，
        每个系统有不同的 API、认证方式和数据格式。
  \item \textbf{幻觉与可靠性}：LLM 的``幻觉''问题
        在企业场景中的代价远高于消费场景。
        消费者可以容忍 ChatGPT 偶尔编造一个不存在的引用；
        但如果企业 AI 助手编造了一个不存在的合同条款、
        一个错误的安全事件、
        或一个虚假的客户信息，
        后果可能是诉讼、声誉损失或安全事故。
\end{itemize}

这些``最后一公里''问题解释了为什么企业 AI 的实际渗透率
始终低于科技公司和分析师的乐观预测——
Microsoft Copilot 的 3.3\% 渗透率就是最好的证明。
但这些挑战也恰恰是创业机会所在：
\textbf{能够系统性地解决这些``最后一公里''问题的公司，
将成为企业 AI 时代的最大赢家。}
Glean 之所以能在 Microsoft Copilot 旁边茁壮成长，
正是因为它比微软更好地解决了权限管理、知识图谱和组织上下文这些
``最后一公里''问题。
\end{warnbox}

\subsection{本章核心公司一览}

为方便读者快速参考，表~\ref{tab:enterprise-companies} 汇总了
本章讨论的主要企业 AI 公司的关键指标。

\begin{table}[htbp]
\centering
\caption{企业 AI 核心公司一览（数据截至 2026 年 3 月）}
\label{tab:enterprise-companies}
\footnotesize
\begin{tabularx}{\textwidth}{lllrX}
\toprule
\rowcolor{figLightGray}
\textbf{公司} & \textbf{赛道} & \textbf{成立} & \textbf{融资/估值} & \textbf{核心指标} \\
\midrule
Microsoft Copilot & 企业 AI 平台 & 2023 & N/A（上市） & 1500 万付费座席，约 54 亿美元 ARR \\
Salesforce & 企业 AI 平台 & 1999 & N/A（上市） & AI ARR 同比增长 114\% \\
ServiceNow & 企业 AI 平台 & 2004 & N/A（上市） & Q4 订阅收入 34.66 亿美元，Now Assist ACV 预计 2026 年达 10 亿 \\
Glean & 企业搜索 & 2019 & 6 亿+ / 72 亿 & ARR 超 1 亿美元，1 亿+ Agent 行为/年 \\
ThoughtSpot & AI 数据分析 & 2012 & 6.74 亿 / 42 亿 & 763 名员工，``Agentic Analytics'' \\
CrowdStrike & AI 网络安全 & 2011 & N/A（上市） & Charlotte AI + AgentWorks，FedRAMP 认证 \\
SentinelOne & AI 网络安全 & 2013 & N/A（上市） & Purple AI 自主安全 \\
Vectra AI & 网络检测 & 2011 & 3.5 亿+ & 2025 Gartner NDR 领导者，35+ AI 专利 \\
o9 Solutions & 供应链 AI & 2009 & 5.36 亿 / 37 亿 & 收入约 2 亿美元，2600 名员工 \\
Altana AI & 供应链合规 & 2018 & 3.24 亿 & 收入 4560 万，美国海关等政府客户 \\
Landing AI & 制造 AI & 2017 & 5700 万+ & 吴恩达创立，ABB Robotics 投资 \\
Buildots & 建筑 AI & 2018 & 1.66 亿 & 三位数收入增长，9 国运营 \\
OpenSpace & 建筑 AI & 2017 & 2 亿 & 收入 1500 万，收购 Disperse.io \\
Avathon & 工业 AI & 2013 & 3 亿 / 14 亿 & 前 BP CEO 任董事长，印度 20\%+ 加油站 \\
UiPath & RPA + AI & 2005 & N/A（上市） & ARR 18.53 亿，GAAP 盈利 \\
Automation Anywhere & RPA + AI & 2003 & 13 亿 / 73 亿 & 收入 4.87 亿，APA 51\% 附着率 \\
来也科技 & RPA + AI & 2015 & 2.11 亿 & 中国唯一 5 年 Gartner MQ \\
\bottomrule
\end{tabularx}
\end{table}

\subsection{企业 AI 的投资启示}

对于投资者和创业者而言，本章的分析提供了几个关键启示：

\textbf{对投资者}：
\begin{itemize}[noitemsep]
  \item 横向平台的投资窗口已经关闭——
        Microsoft、Google、Salesforce、ServiceNow 四巨头已经占位，
        新的横向平台创业公司几乎不可能突围。
        但``右下象限''（纵向深层）仍有大量未被占据的赛道——
        每个行业都可能诞生一个 Glean 级别的垂直 AI 冠军。
  \item 关注``数据飞轮速度''而非``当前收入''。
        企业 AI 公司的长期价值取决于其数据积累速度——
        Glean 在 1000 家企业中构建的知识图谱、
        CrowdStrike 在数百万端点上积累的威胁情报、
        Buildots 在数千个建筑项目中采集的影像数据——
        这些数据资产的价值随时间指数级增长。
  \item 并购是主要的退出路径。
        独立 IPO 对大多数垂直 AI 公司来说不现实
        （市场规模和增速不足以支撑独立上市的估值期望），
        被平台巨头收购（如 Moveworks 被 ServiceNow 以 28.5 亿美元收购）
        是更现实的退出场景。
\end{itemize}

\textbf{对创业者}：
\begin{itemize}[noitemsep]
  \item 选择``右下象限''（纵向深层）作为起点。
        不要试图做``下一个 Microsoft Copilot''，
        而是成为某个行业的``不可替代的 AI 基础设施''。
  \item ``从搜索/可见性切入，向 Agent 扩展''是被验证的增长路径。
        Glean 从搜索到 Agent，ThoughtSpot 从查询到 Agentic Analytics，
        Buildots 从影像监控到预测性分析——
        都遵循了这一``先解决信息获取问题，
        再解决决策和执行问题''的路径。
  \item 安全和权限不是``事后补丁''，而是``第一天的架构决策''。
        Glean 的权限感知设计、CrowdStrike 的 ISO 42001 认证、
        来也科技的信创适配——
        这些安全和合规能力是企业客户决策的先决条件，
        不是产品成熟后再添加的功能。
  \item 定价模型是竞争力的核心。
        Salesforce 从``按座席''到``按对话''的定价转型
        直接推动了 114\% 的 AI ARR 增长。
        如何将 AI 的价值准确地映射到客户愿意支付的价格上，
        是企业 AI 创业的核心课题之一。
\end{itemize}

本章的最终结论是：
\textbf{企业 AI 不是一个赢家通吃的市场，而是一个``层层渗透''的生态系统}。
横向平台（Microsoft、Salesforce、ServiceNow）定义了基础设施层，
横向深层公司（Glean、ThoughtSpot）占据了跨行业的关键功能节点，
纵向深层公司（CrowdStrike、Buildots、Landing AI）
在各自的行业中建立了不可替代的数据壁垒。
这三个层次的公司不是简单的竞争关系，
而是构成了一个相互依存、相互竞争、不断演化的复杂生态。
理解这个生态的结构和演化方向，
是把握企业 AI 创业和投资机会的关键。
